% ==============================================================================
% APÊNDICES
% Deve ser incluído no documento principal APÓS \appendix
% Exemplo de uso no main.tex:
%   \appendix
%   % ==============================================================================
% APÊNDICES
% Deve ser incluído no documento principal APÓS \appendix
% Exemplo de uso no main.tex:
%   \appendix
%   % ==============================================================================
% APÊNDICES
% Deve ser incluído no documento principal APÓS \appendix
% Exemplo de uso no main.tex:
%   \appendix
%   % ==============================================================================
% APÊNDICES
% Deve ser incluído no documento principal APÓS \appendix
% Exemplo de uso no main.tex:
%   \appendix
%   \input{Apendices}
%
% Pacotes necessários no preâmbulo do documento principal:
%   \usepackage{listings}
%   \usepackage{xcolor}
% ==============================================================================

% Configuração global das listagens de código Python
\lstset{
  language=Python,
  basicstyle=\ttfamily\footnotesize,
  keywordstyle=\color{blue}\bfseries,
  commentstyle=\color{gray}\itshape,
  stringstyle=\color{teal},
  numberstyle=\tiny\color{gray},
  numbers=left,
  stepnumber=1,
  numbersep=8pt,
  backgroundcolor=\color{gray!6},
  frame=single,
  framerule=0.4pt,
  rulecolor=\color{gray!40},
  breaklines=true,
  breakatwhitespace=false,
  showstringspaces=false,
  tabsize=4,
  captionpos=b,
  morekeywords={self, True, False, None, yield, lambda, with, as, assert,
                del, elif, except, finally, from, global, import, in,
                is, not, or, pass, raise, try},
  literate={á}{{\'a}}1 {é}{{\'e}}1 {í}{{\'i}}1 {ó}{{\'o}}1 {ú}{{\'u}}1
           {â}{{\^a}}1 {ê}{{\^e}}1 {î}{{\^i}}1 {ô}{{\^o}}1 {û}{{\^u}}1
           {ã}{{\~a}}1 {õ}{{\~o}}1 {ç}{{\c{c}}}1
           {Á}{{\'A}}1 {É}{{\'E}}1 {Í}{{\'I}}1 {Ó}{{\'O}}1 {Ú}{{\'U}}1
           {Â}{{\^A}}1 {Ê}{{\^E}}1 {Ô}{{\^O}}1 {Ã}{{\~A}}1 {Õ}{{\~O}}1
           {Ç}{{\c{C}}}1,
}

\lstdefinestyle{bash}{
  language=bash,
  keywordstyle=\color{violet}\bfseries,
  commentstyle=\color{gray}\itshape,
  stringstyle=\color{teal},
  backgroundcolor=\color{gray!8},
}

\chapter{Código-fonte da \textit{Engine} Matemática (\texttt{engine.py})}
\label{ap:engine}

Este apêndice apresenta o código-fonte completo dos componentes fundamentais da \textit{engine} matemática do \textit{Hephaestus}, implementados no módulo \texttt{engine.py}. As listagens cobrem as funções de nível de módulo utilizadas para paralelismo, a classe \texttt{RobotMathEngine} (ambiente seco) e a subclasse \texttt{RobotMathHydro} (ambiente subaquático). O código é apresentado na íntegra para viabilizar reprodutibilidade e auditoria matemática das equações derivadas.

% ------------------------------------------------------------------------------
\section*{A.1 Funções de Nível de Módulo para Paralelismo}
% ------------------------------------------------------------------------------

As duas funções a seguir são definidas no escopo do módulo para permitir sua serialização via \texttt{pickle}, requisito do \texttt{ProcessPoolExecutor} no Python. A função \texttt{\_diff\_M\_col} diferencia a matriz de inércia em relação a uma variável de junta $q_k$; a função \texttt{\_calc\_coriolis\_row} monta a $i$-ésima linha do vetor de Coriolis pela fórmula dos símbolos de Christoffel.

\begin{lstlisting}[language=Python, caption={Funções paralelas para derivação de $\partial M/\partial q_k$ e cálculo de $C_i$.}, label={lst:parallel_funcs}]
from concurrent.futures import ProcessPoolExecutor
from itertools import repeat
import os, time
import sympy as sp
from sympy.physics.mechanics import dynamicsymbols

def _diff_M_col(args):
    """Diferencia a matriz M em relação a uma variável simbólica qk.
    Executado em processo separado para paralelizar o cálculo de dM/dq.
    """
    M_list, n, qk = args
    M = sp.Matrix(n, n, M_list)
    return list(M.diff(qk))


def _calc_coriolis_row(i, q, dq, dM_dq):
    n = len(q)
    termo_linha = sp.Integer(0)
    for j in range(n):
        for k in range(j, n):
            dM_ij_dk = dM_dq[k][i, j]
            dM_ik_dj = dM_dq[j][i, k]
            dM_jk_di = dM_dq[i][j, k]
            c_ijk = sp.Rational(1, 2) * (dM_ij_dk + dM_ik_dj - dM_jk_di)
            if c_ijk != 0:
                termo = c_ijk * dq[j] * dq[k]
                if k != j:
                    termo *= 2
                termo_linha += termo
    # sp.collect() omitido intencionalmente: custo superlinear sem benefício,
    # pois o cse=True no lambdify posterior realiza compactação sistematica.
    return termo_linha
\end{lstlisting}

% ------------------------------------------------------------------------------
\section*{A.2 Classe \texttt{RobotMathEngine} — Ambiente Seco}
% ------------------------------------------------------------------------------

\begin{lstlisting}[language=Python, caption={Construtor e configuração de símbolos de \texttt{RobotMathEngine}.}, label={lst:engine_init}]
class RobotMathEngine:
    def __init__(self, joint_config, link_vectors_mask,
                 logger_callback=None, num_workers=None):
        self.joint_config    = joint_config
        self.link_vectors_mask = [sp.Matrix(v) for v in link_vectors_mask]
        self.log             = logger_callback if logger_callback else print
        self.num_workers     = num_workers

        self.t = sp.symbols('t')
        self.g = sp.symbols('g')

        self.q, self.dq, self.params_list   = [], [], []
        self.frames, self.rotation_matrices  = [], []
        self.com_positions_global            = []
        self.angular_velocities              = []
        self.masses                          = []

        self.params_list.append(self.g)   # g sempre presente como parâmetro
        self.M, self.G_vec, self.C_total, self.Jacobian = None, None, None, None

    def _rot_matrix_local(self, axis, angle):
        c, s = sp.cos(angle), sp.sin(angle)
        if axis == 'x': return sp.Matrix([[1,0,0],[0,c,-s],[0,s,c]])
        if axis == 'y': return sp.Matrix([[c,0,s],[0,1,0],[-s,0,c]])
        if axis == 'z': return sp.Matrix([[c,-s,0],[s,c,0],[0,0,1]])
        return sp.eye(3)

    def _get_axis_vector(self, axis):
        if axis == 'x': return sp.Matrix([1, 0, 0])
        if axis == 'y': return sp.Matrix([0, 1, 0])
        if axis == 'z': return sp.Matrix([0, 0, 1])
        return sp.Matrix([0, 0, 0])
\end{lstlisting}

\begin{lstlisting}[language=Python, caption={Cinemática simbólica direta: transformações homogêneas e posições de CM.}, label={lst:step1}]
    def step_1_kinematics(self):
        self.log("1. (AR) Calculando Cinemarica (Com CM Arbitrario)...")
        T_acc   = sp.eye(4)
        R_acc   = sp.eye(3)
        omega_acc = sp.Matrix([0, 0, 0])

        for i, (j_type, link_vec) in enumerate(
                zip(self.joint_config, self.link_vectors_mask)):
            q  = dynamicsymbols(f'q{i+1}')
            dq = dynamicsymbols(f'q{i+1}', 1)
            self.q.append(q);  self.dq.append(dq)

            m  = sp.symbols(f'm{i+1}')
            L  = sp.symbols(f'L{i+1}')
            cx, cy, cz = sp.symbols(f'cx{i+1} cy{i+1} cz{i+1}')

            self.masses.append(m)
            self.params_list.extend([m, L, cx, cy, cz])

            type_char = j_type[0]
            axis_char = j_type[1].lower()

            if type_char == 'R':            # Junta de revolucao
                axis_vec_global = R_acc * self._get_axis_vector(axis_char)
                omega_new = omega_acc + axis_vec_global * dq
                R_j = self._rot_matrix_local(axis_char, q)
                P_j = sp.Matrix([0, 0, 0])
            elif type_char == 'D':          # Junta prismatica
                omega_new = omega_acc
                R_j = sp.eye(3)
                P_j = sp.Matrix([0, 0, 0])
                if axis_char == 'x': P_j[0] = q
                if axis_char == 'y': P_j[1] = q
                if axis_char == 'z': P_j[2] = q

            R_acc = R_acc * R_j
            self.rotation_matrices.append(R_acc)
            self.angular_velocities.append(omega_new)
            omega_acc = omega_new

            T_joint = sp.eye(4)
            T_joint[0:3, 0:3] = R_j
            T_joint[0:3, 3]   = P_j
            T_at_joint_start  = T_acc * T_joint

            P_link_vec        = link_vec * L
            T_link            = sp.eye(4)
            T_link[0:3, 3]    = P_link_vec
            T_acc             = T_at_joint_start * T_link
            self.frames.append(T_acc)

            # CM Global: transforma coordenadas locais pelo frame ate a junta
            v_cm_local   = sp.Matrix([cx, cy, cz, 1])
            p_cm_global  = T_at_joint_start * v_cm_local
            self.com_positions_global.append(p_cm_global[0:3, 0])
\end{lstlisting}

\begin{lstlisting}[language=Python, caption={Derivação simbólica de $M(q)$, $G(q)$ e $J(q)$ pelo método do Jacobiano da inércia.}, label={lst:step2}]
    def step_2_jacobian_M_G(self):
        self.log("2. (AR) Dinamica M e G (Steiner + Tensores)...")
        _t0 = time.perf_counter()
        n   = len(self.q)
        self.M = sp.zeros(n, n)
        V_tot  = 0

        for i in range(n):
            m       = self.masses[i]
            pos_cm  = self.com_positions_global[i]
            R_global = self.rotation_matrices[i]

            J_v = pos_cm.jacobian(self.q)         # Jacobiano linear do CM
            J_w = self.angular_velocities[i].jacobian(self.dq)  # Jacobiano angular

            Ixx, Iyy, Izz = sp.symbols(f'Ixx{i+1} Iyy{i+1} Izz{i+1}')
            self.params_list.extend([Ixx, Iyy, Izz])

            # Transformacao de Steiner: I_global = R * I_local * R^T
            I_local  = sp.Matrix([[Ixx, 0, 0], [0, Iyy, 0], [0, 0, Izz]])
            I_global = R_global * I_local * R_global.T

            # M += m*Jv^T*Jv + Jw^T * I_global * Jw  (metodo do Jacobiano)
            self.M  += m * J_v.T * J_v + J_w.T * I_global * J_w

            V_tot   += m * self.g * pos_cm[2]   # Energia potencial gravitacional
            self.log(f"  M/G: elo {i+1}/{n} concluido")

        self.G_vec   = sp.Matrix([V_tot]).jacobian(self.q).T
        J_lin        = self.frames[-1][0:3, 3].jacobian(self.q)
        J_ang        = self.angular_velocities[-1].jacobian(self.dq)
        self.Jacobian = J_lin.col_join(J_ang)    # Jacobiano 6xN completo
        self.log(f"  M/G/J: concluido em {time.perf_counter()-_t0:.1f}s")
\end{lstlisting}

\begin{lstlisting}[language=Python, caption={Cálculo do vetor de Coriolis pelos símbolos de Christoffel com paralelismo opcional.}, label={lst:step3}]
    def step_3_coriolis_combined(self):
        self.log("3. Calculando Coriolis...")
        n          = len(self.q)
        self.C_total = sp.zeros(n, 1)
        cpu_total  = os.cpu_count() or 1
        workers    = self.num_workers if self.num_workers is not None else cpu_total
        use_parallel = workers is not None and workers > 1

        # --- Passo 3a: dM/dq  (paralelizavel por coluna) ---
        self.log("  dM/dq: diferenciando M" +
                 (" em paralelo" if use_parallel else "") + "...")
        _t0 = time.perf_counter()
        M_list = list(self.M)
        tasks  = [(M_list, n, qk) for qk in self.q]

        if use_parallel:
            with ProcessPoolExecutor(max_workers=workers) as ex:
                dM_dq_flat = list(ex.map(_diff_M_col, tasks))
        else:
            dM_dq_flat = [_diff_M_col(t) for t in tasks]

        dM_dq = [sp.Matrix(n, n, flat) for flat in dM_dq_flat]
        self.log(f"  dM/dq: concluido em {time.perf_counter()-_t0:.1f}s")

        # --- Passo 3b: Simbolos de Christoffel  C_i(q,dq) ---
        if use_parallel:
            with ProcessPoolExecutor(max_workers=workers) as ex:
                for i, row in enumerate(
                    ex.map(_calc_coriolis_row,
                           range(n), repeat(self.q),
                           repeat(self.dq), repeat(dM_dq)),
                    start=1):
                    self.C_total[i - 1] = row
                    self.log(f"  Coriolis: linha {i}/{n} concluida")
        else:
            for i in range(n):
                self.C_total[i] = _calc_coriolis_row(i, self.q, self.dq, dM_dq)
                self.log(f"  Coriolis: linha {i+1}/{n} concluida")
\end{lstlisting}

\begin{lstlisting}[language=Python, caption={Substituição de símbolos dinâmicos e exportação do dicionário de modelo.}, label={lst:step4}]
    def step_4_prepare_export(self):
        self.log("4. Otimizando equacoes...")
        mapa_subs = {}
        for i in range(len(self.q)):
            mapa_subs[self.q[i]]  = sp.Symbol(f'q{i+1}')
            mapa_subs[self.dq[i]] = sp.Symbol(f'dq{i+1}')

        dq_vec      = sp.Matrix(self.dq)
        V_cartesian = self.Jacobian * dq_vec if self.Jacobian is not None else None

        return {
            "M":      self.M,
            "G":      self.G_vec,
            "C":      self.C_total,
            "J":      self.Jacobian,
            "FK_Pos": self.frames[-1],
            "FK_Vel": V_cartesian,
            "Subs":   mapa_subs,
            "Mode":   "Air"
        }
\end{lstlisting}

% ------------------------------------------------------------------------------
\section*{A.3 Subclasse \texttt{RobotMathHydro} — Extensão Hidrodinâmica}
% ------------------------------------------------------------------------------

\begin{lstlisting}[language=Python, caption={Extensão hidrodinâmica: massa adicionada, empuxo e potencial aparente.}, label={lst:hydro}]
class RobotMathHydro(RobotMathEngine):
    def __init__(self, joint_config, link_vectors_mask,
                 logger_callback=None, num_workers=None):
        super().__init__(joint_config, link_vectors_mask,
                         logger_callback, num_workers)
        self.rho     = sp.symbols('rho')   # densidade do fluido
        self.volumes = []

    def step_1_kinematics_hydro(self):
        super().step_1_kinematics()
        self.log("   -> Adicionando parametros hidrodinamicos...")
        for i in range(len(self.q)):
            vol = sp.symbols(f'vol{i+1}')
            self.volumes.append(vol)
            self.params_list.append(vol)   # vol_i adicionado a assinatura

    def step_2_jacobian_M_G(self):
        self.log("2. (AGUA) Dinamica: M (Inercia + Added Mass) e G (Peso - Empuxo)...")
        _t0  = time.perf_counter()
        n    = len(self.q)
        self.M = sp.zeros(n, n)
        V_tot  = 0

        for i in range(n):
            m       = self.masses[i]
            vol     = self.volumes[i]
            pos_cm  = self.com_positions_global[i]
            R_global = self.rotation_matrices[i]

            J_v = pos_cm.jacobian(self.q)
            J_w = self.angular_velocities[i].jacobian(self.dq)

            # Tensor de inercia rigido (referencial global via Steiner)
            Ixx, Iyy, Izz = sp.symbols(f'Ixx{i+1} Iyy{i+1} Izz{i+1}')
            self.params_list.extend([Ixx, Iyy, Izz])
            I_local_RB  = sp.Matrix([[Ixx,0,0],[0,Iyy,0],[0,0,Izz]])
            I_global_RB = R_global * I_local_RB * R_global.T

            # Tensor de massa adicionada (diagonal local, 6 componentes por elo)
            ma_u, ma_v, ma_w = sp.symbols(f'ma_u{i+1} ma_v{i+1} ma_w{i+1}')
            ma_p, ma_q, ma_r = sp.symbols(f'ma_p{i+1} ma_q{i+1} ma_r{i+1}')
            self.params_list.extend([ma_u, ma_v, ma_w, ma_p, ma_q, ma_r])

            MA_lin_local = sp.Matrix([[ma_u,0,0],[0,ma_v,0],[0,0,ma_w]])
            MA_rot_local = sp.Matrix([[ma_p,0,0],[0,ma_q,0],[0,0,ma_r]])

            # Rotacao para referencial global: M_A^(g) = R * M_A^(l) * R^T
            MA_lin_global = R_global * MA_lin_local * R_global.T
            MA_rot_global = R_global * MA_rot_local * R_global.T

            # M_hydro += Jv^T (m*I + M_A,lin)^(g) Jv + Jw^T (I_RB + M_A,rot)^(g) Jw
            self.M += J_v.T * (m * sp.eye(3) + MA_lin_global) * J_v
            self.M += J_w.T * (I_global_RB + MA_rot_global) * J_w

            # Potencial aparente: V_i = (m - rho*vol) * g * z_cm  (Peso - Empuxo)
            peso_aparente = (m - self.rho * vol) * self.g
            V_tot += peso_aparente * pos_cm[2]
            self.log(f"  M/G: elo {i+1}/{n} concluido")

        self.G_vec   = sp.Matrix([V_tot]).jacobian(self.q).T
        J_lin        = self.frames[-1][0:3, 3].jacobian(self.q)
        J_ang        = self.angular_velocities[-1].jacobian(self.dq)
        self.Jacobian = J_lin.col_join(J_ang)
        self.log(f"  M/G/J: concluido em {time.perf_counter()-_t0:.1f}s")

    def step_4_prepare_export(self):
        data = super().step_4_prepare_export()
        data["Mode"] = "Hydro"
        return data
\end{lstlisting}

% ==============================================================================
\chapter{Código-fonte dos Controladores e do Simulador (\texttt{simulator.py})}
\label{ap:simulador}

Este apêndice apresenta as classes de controle e a estrutura principal da classe \texttt{RobotSimulator} implementadas no módulo \texttt{simulator.py}. As listagens incluem os controladores CT-SMC, Super-Twisting (STA) e LADRC descentralizado, bem como o construtor do simulador, que realiza a compilação das expressões simbólicas via \texttt{lambdify} com CSE, e a infraestrutura de avaliação paralela numérica.

% ------------------------------------------------------------------------------
\section*{B.1 Funções Auxiliares de Orientação}
% ------------------------------------------------------------------------------

\begin{lstlisting}[language=Python, caption={Erro de orientação pelo operador de anel (skew-symmetric) e SLERP em $SO(3)$.}, label={lst:orient_helpers}]
import numpy as np

def _rotation_error_vec(R_ref, R_curr):
    """Erro de orientacao 3D pela formulacao de anel (skew-symmetric).

    Retorna e_R = vee(0.5*(R_ref @ R_curr^T - R_curr @ R_ref^T)),
    proporcional ao eixo-angulo para erros pequenos.
    """
    skew = 0.5 * (R_ref @ R_curr.T - R_curr @ R_ref.T)
    return np.array([skew[2, 1], skew[0, 2], skew[1, 0]])


def _slerp_rotation(R0, Rf, s, sd):
    """SLERP entre R0 e Rf na fracao s em [0, 1].

    Parâmetros
    ----------
    R0, Rf : matrizes de rotacao 3x3
    s      : fracao de interpolacao atual  (0 = R0, 1 = Rf)
    sd     : derivada temporal de s (ds/dt)

    Retorna
    -------
    R_ref     : rotacao interpolada 3x3
    omega_ref : velocidade angular no referencial global (rad/s)
    """
    dR        = R0.T @ Rf
    trace_val = np.clip((np.trace(dR) - 1.0) / 2.0, -1.0, 1.0)
    theta_total = np.arccos(trace_val)
    if abs(theta_total) < 1e-9:
        return R0.copy(), np.zeros(3)
    sin_t      = np.sin(theta_total)
    skew       = (dR - dR.T) / (2.0 * sin_t)
    axis_local = np.array([skew[2, 1], skew[0, 2], skew[1, 0]])
    theta_s    = s * theta_total
    K = np.array([[ 0.0,           -axis_local[2],  axis_local[1]],
                  [ axis_local[2],  0.0,            -axis_local[0]],
                  [-axis_local[1],  axis_local[0],   0.0          ]])
    dR_s  = np.eye(3) + np.sin(theta_s) * K + (1.0 - np.cos(theta_s)) * (K @ K)
    R_ref = R0 @ dR_s
    omega_ref = (sd * theta_total) * (R0 @ axis_local)
    return R_ref, omega_ref
\end{lstlisting}

% ------------------------------------------------------------------------------
\section*{B.2 Controlador CT-SMC com Camada Limite}
% ------------------------------------------------------------------------------

\begin{lstlisting}[language=Python, caption={Controle por Torque Computado com Modo Deslizante (CT-SMC).}, label={lst:smc}]
class SlidingModeControl:
    """Controlador CT-SMC: Computed Torque + Sliding Mode com camada limite.

    Lei de controle:
        q_ddot_f = alpha*q_ddot_d + (1-alpha)*q_ddot_f  (filtro passabaixa)
        v        = q_ddot_f - lambda*e_dot - K*tanh(S/phi)
        tau      = M(q)*v + C(q,dq) + G(q)
    onde S = e_dot + lambda*e  (superficie deslizante).

    Referencia: Slotine & Li, Applied Nonlinear Control, 1991.
    """

    def __init__(self, num_dof, lambda_gain, K, phi, ddq_filter_alpha=1.0):
        self.lambda_gain    = (np.full(num_dof, lambda_gain, dtype=float)
                               if np.isscalar(lambda_gain)
                               else np.array(lambda_gain, dtype=float))
        self.K              = (np.full(num_dof, K, dtype=float)
                               if np.isscalar(K) else np.array(K, dtype=float))
        self.phi            = (np.full(num_dof, phi, dtype=float)
                               if np.isscalar(phi) else np.array(phi, dtype=float))
        self.phi            = np.where(self.phi == 0, 1e-6, self.phi)
        self.alpha          = float(np.clip(ddq_filter_alpha, 1e-6, 1.0))
        self._ddq_d_filt    = None

    def compute_tau(self, q, dq, q_d, dq_d, ddq_d, M, C, G, dt=0.0):
        if self._ddq_d_filt is None:
            self._ddq_d_filt = ddq_d.copy()
        else:
            self._ddq_d_filt = (self.alpha * ddq_d
                                + (1.0 - self.alpha) * self._ddq_d_filt)
        e       = q - q_d
        e_dot   = dq - dq_d
        S       = e_dot + self.lambda_gain * e
        v       = (self._ddq_d_filt - self.lambda_gain * e_dot
                   - self.K * np.tanh(S / self.phi))
        return M @ v + C + G
\end{lstlisting}

% ------------------------------------------------------------------------------
\section*{B.3 Controlador Super-Twisting (STA)}
% ------------------------------------------------------------------------------

\begin{lstlisting}[language=Python, caption={Controlador Super-Twisting de 2ª ordem (STA) com sinal de controle contínuo.}, label={lst:sta}]
class SuperTwistingSMC:
    """Controlador Super-Twisting SMC (STA) baseado em modelo.

    Algoritmo de modo deslizante de 2a ordem:
        q_ddot_f = alpha*q_ddot_d + (1-alpha)*q_ddot_f
        v_sw     = -k1*|S|^(1/2)*sign(S) + z
        z_dot    = -k2*sign(S)
        tau      = M(q)*[q_ddot_f - lambda*e_dot + v_sw] + C(q,dq) + G(q)

    Propriedades: convergencia em tempo finito, sinal de controle continuo.
    Condicoes de ganho (Moreno & Osorio, 2012): k2 > delta, k1 > 2*sqrt(k2*delta).

    Referencias:
    - Levant, Int. J. Control, 58(6), 1993.
    - Moreno & Osorio, IEEE Trans. Autom. Control, 57(4), 2012.
    """

    def __init__(self, num_dof, lambda_gain, k1, k2, ddq_filter_alpha=1.0):
        self.lambda_gain = (np.full(num_dof, lambda_gain, dtype=float)
                            if np.isscalar(lambda_gain)
                            else np.array(lambda_gain, dtype=float))
        self.k1          = (np.full(num_dof, k1, dtype=float)
                            if np.isscalar(k1) else np.array(k1, dtype=float))
        self.k2          = (np.full(num_dof, k2, dtype=float)
                            if np.isscalar(k2) else np.array(k2, dtype=float))
        self.z           = np.zeros(num_dof)
        self.alpha       = float(np.clip(ddq_filter_alpha, 1e-6, 1.0))
        self._ddq_d_filt = None

    def reset(self):
        self.z[:]        = 0.0
        self._ddq_d_filt = None

    def compute_tau(self, q, dq, q_d, dq_d, ddq_d, M, C, G, dt):
        if self._ddq_d_filt is None:
            self._ddq_d_filt = ddq_d.copy()
        else:
            self._ddq_d_filt = (self.alpha * ddq_d
                                + (1.0 - self.alpha) * self._ddq_d_filt)
        e     = q - q_d
        e_dot = dq - dq_d
        S     = e_dot + self.lambda_gain * e

        # Parcela de chaveamento continua do STA
        v_sw = -self.k1 * np.abs(S) ** 0.5 * np.sign(S) + self.z

        v        = self._ddq_d_filt - self.lambda_gain * e_dot + v_sw
        self.z  -= self.k2 * np.sign(S) * dt    # integracao de z (Euler explicito)
        return M @ v + C + G
\end{lstlisting}

% ------------------------------------------------------------------------------
\section*{B.4 LADRC Descentralizado com ESO de Terceira Ordem}
% ------------------------------------------------------------------------------

\begin{lstlisting}[language=Python, caption={LADRC descentralizado: ESO de 3ª ordem e lei de controle por junta.}, label={lst:ladrc}]
class Decentralized_LADRC:
    """LADRC descentralizado para multiplas juntas (2a ordem por junta).

    Para a junta i: ddot_q_i = b0_i * tau_i + f_i  (f_i = disturbio total).
    O ESO de 3a ordem (linear) estima [q_i, dq_i, f_i] em z1, z2, z3.
    Parametrizacao dos ganhos (Gao, 2003): beta1=3*wo, beta2=3*wo^2, beta3=wo^3.
    """

    def __init__(self, num_dof, b0, wo, kp, kd, dt,
                 z_limit=None, tau_limit=None,
                 max_wo_dt=0.1, z3_filter_alpha=1.0):
        self.b0  = (np.full(num_dof, b0, dtype=float)
                    if np.isscalar(b0) else np.array(b0, dtype=float))
        self.wo  = (np.full(num_dof, wo, dtype=float)
                    if np.isscalar(wo) else np.array(wo, dtype=float))
        self.kp  = (np.full(num_dof, kp, dtype=float)
                    if np.isscalar(kp) else np.array(kp, dtype=float))
        self.kd  = (np.full(num_dof, kd, dtype=float)
                    if np.isscalar(kd) else np.array(kd, dtype=float))
        self.dt  = dt
        self.z_limit        = z_limit
        self.tau_limit      = tau_limit
        self.max_wo_dt      = max_wo_dt
        self.z3_filter_alpha = float(np.clip(z3_filter_alpha, 0.0, 1.0))
        self.z3_filtered    = None

        self._refresh_gains()
        self.z1 = np.zeros(num_dof, dtype=float)
        self.z2 = np.zeros(num_dof, dtype=float)
        self.z3 = np.zeros(num_dof, dtype=float)

    def _refresh_gains(self):
        # Limita wo para garantir estabilidade do ESO discreto: wo*dt <= max_wo_dt
        if self.max_wo_dt is not None:
            wo_dt = self.wo * self.dt
            self.wo = np.where(wo_dt > self.max_wo_dt,
                               self.max_wo_dt / self.dt, self.wo)
        self.beta1 = 3.0 *  self.wo
        self.beta2 = 3.0 * (self.wo ** 2)
        self.beta3 =         self.wo ** 3

    def reset_state(self, q, dq, z3=None):
        self.z1 = np.array(q,  dtype=float).copy()
        self.z2 = np.array(dq, dtype=float).copy()
        self.z3 = (np.zeros(self.num_dof, dtype=float) if z3 is None
                   else (np.full(self.num_dof, z3, dtype=float)
                         if np.isscalar(z3) else np.array(z3, dtype=float)))
        self.z3_filtered = self.z3.copy()

    def update_eso(self, q, tau_prev):
        """Avanca o ESO por um passo de integracao (Euler explicito)."""
        error  = self.z1 - q
        z1_dot = self.z2 - self.beta1 * error
        z2_dot = self.z3 - self.beta2 * error + self.b0 * tau_prev
        z3_dot =         - self.beta3 * error

        self.z1 += z1_dot * self.dt
        self.z2 += z2_dot * self.dt
        self.z3 += z3_dot * self.dt

        # Filtro IIR de 1a ordem em z3 (atenua alta frequencia antes da lei)
        a = self.z3_filter_alpha
        self.z3_filtered = a * self.z3 + (1.0 - a) * self.z3_filtered

    def compute_control(self, q_d, dq_d, ddq_d, q_meas=None, dq_meas=None):
        """Calcula o torque de controle LADRC (cancela disturbio + PD + ff)."""
        q_fb  = q_meas  if q_meas  is not None else self.z1
        dq_fb = dq_meas if dq_meas is not None else self.z2

        z3_use = self.z3_filtered if self.z3_filtered is not None else self.z3
        u0     = ddq_d + self.kp * (q_d - q_fb) + self.kd * (dq_d - dq_fb)
        b0_safe = np.where(self.b0 == 0.0, 1e-6, self.b0)
        u = (u0 - z3_use) / b0_safe
        if self.tau_limit is not None:
            u = np.clip(u, -self.tau_limit, self.tau_limit)
        return u
\end{lstlisting}

% ------------------------------------------------------------------------------
\section*{B.5 Compilação Simbólico-Numérica e Executor Persistente}
% ------------------------------------------------------------------------------

\begin{lstlisting}[language=Python, caption={Funções de worker para avaliação paralela de $M$, $C$, $G$ no executor persistente.}, label={lst:workers}]
import sympy as sp

_WORKER_FUNCS = {}   # dicionario de modulo: preenchido por _init_worker


def _init_worker(sym_vars, expr_M, expr_C, expr_G):
    """Inicializador do worker: compila lambdify localmente no processo."""
    global _WORKER_FUNCS
    _WORKER_FUNCS = {
        "M": sp.lambdify(sym_vars, expr_M, modules="numpy"),
        "C": sp.lambdify(sym_vars, expr_C, modules="numpy"),
        "G": sp.lambdify(sym_vars, expr_G, modules="numpy"),
    }


def _eval_worker(task):
    """Avalia uma das funcoes compiladas com os argumentos numericos fornecidos."""
    func_name, args = task
    return _WORKER_FUNCS[func_name](*args)
\end{lstlisting}

\begin{lstlisting}[language=Python, caption={Construtor de \texttt{RobotSimulator}: compilação \texttt{lambdify} + CSE e inicialização do executor.}, label={lst:sim_init}]
from concurrent.futures import ProcessPoolExecutor
import time, numpy as np, sympy as sp

class RobotSimulator:
    def __init__(self, robot_math_instance, mode="Air"):
        self.bot       = robot_math_instance
        self.mode      = mode
        self.num_dof   = len(self.bot.q)
        self.q_home    = np.zeros(self.num_dof)
        self._executor = None      # executor persistente (inicializado apos compilacao)
        self._executor_workers = 0

        print(f"[{mode}] Compilando equacoes com lambdify+CSE...")
        _t0 = time.perf_counter()

        # Identificacao de juntas rotacionais (para wrap em [-pi, pi])
        self.is_rotational = np.array(
            [j.startswith('R') for j in self.bot.joint_config], dtype=bool)

        # Assinatura completa: [q1..qN, dq1..dqN, params...]
        self.sym_vars = self.bot.q + self.bot.dq + self.bot.params_list
        if hasattr(self.bot, 'rho'):
            self.sym_vars.append(self.bot.rho)

        # Compilacao separada por grandeza (M, C, G, J, FK)
        # cse=True: CSE antes da geracao do codigo -> reducao de 60-90% no bytecode
        self.expr_M = self.bot.M
        self.expr_C = self.bot.C_total
        self.expr_G = self.bot.G_vec

        self.func_M = sp.lambdify(self.sym_vars, self.expr_M,
                                  modules="numpy", cse=True)
        self.func_C = sp.lambdify(self.sym_vars, self.expr_C,
                                  modules="numpy", cse=True)
        self.func_G = sp.lambdify(self.sym_vars, self.expr_G,
                                  modules="numpy", cse=True)
        self.func_J = sp.lambdify(self.sym_vars, self.bot.Jacobian,
                                  modules="numpy", cse=True)

        # Funcoes FK por elo para visualizacao animada
        self.funcs_fk_all_links = [
            sp.lambdify(self.sym_vars,
                        frame[0:3, 3].subs(self.bot.step_4_prepare_export()["Subs"]),
                        modules="numpy", cse=True)
            for frame in self.bot.frames
        ]
        print(f"Compilacao concluida em {time.perf_counter()-_t0:.1f}s")

    def warm_up_workers(self, num_workers):
        """Inicializa o executor persistente e aquece os workers."""
        if (self._executor is not None
                and self._executor_workers == num_workers):
            return   # ja inicializado com mesmo numero de workers
        if self._executor is not None:
            self._executor.shutdown(wait=False)
        self._executor = ProcessPoolExecutor(
            max_workers=num_workers,
            initializer=_init_worker,
            initargs=(self.sym_vars, self.expr_M, self.expr_C, self.expr_G))
        self._executor_workers = num_workers
        # Warmup: submete tarefas dummy para pre-compilar os workers
        dummy_args = tuple(np.zeros(len(self.sym_vars)))
        futs = [self._executor.submit(_eval_worker, (k, dummy_args))
                for k in ("M", "C", "G")]
        for f in futs:
            try: f.result()
            except Exception: pass

    def close(self):
        """Encerra o executor persistente de forma limpa."""
        if self._executor is not None:
            self._executor.shutdown(wait=False)
            self._executor = None
\end{lstlisting}

% ==============================================================================
\chapter{Requisitos de Software, Instalação e Execução do \textit{Hephaestus}}
\label{ap:instalacao}

Este apêndice descreve os requisitos de software, o procedimento de instalação e as instruções de execução do ambiente \textit{Hephaestus} a partir do código-fonte.

% ------------------------------------------------------------------------------
\section*{C.1 Requisitos de Sistema}
% ------------------------------------------------------------------------------

O \textit{Hephaestus} foi desenvolvido e validado no seguinte ambiente:

\begin{itemize}
    \item \textbf{Sistema Operacional:} Windows 10/11 (64 bits). O código é multiplataforma e compatível com Linux e macOS com as mesmas dependências; no entanto, a interface gráfica com \textit{CustomTkinter} pode apresentar diferenças visuais menores entre plataformas.
    \item \textbf{Python:} versão 3.10 ou superior (recomendado: 3.11). O uso de \texttt{ProcessPoolExecutor} com método de \textit{spawn} requer Python $\geq$ 3.8; a tipagem e as dependências utilizadas requerem Python $\geq$ 3.10.
    \item \textbf{Hardware recomendado:} processador com pelo menos 4 núcleos lógicos (para aproveitar o paralelismo de derivação); 8 GB de RAM (para modelos com $N \geq 6$ GDL, cujas matrizes simbólicas podem ocupar centenas de MB antes da compilação).
\end{itemize}

% ------------------------------------------------------------------------------
\section*{C.2 Dependências Python}
% ------------------------------------------------------------------------------

A Tabela~\ref{tab:deps} lista as dependências diretas do projeto com as versões mínimas recomendadas.

\begin{table}[h!]
\centering
\caption{Dependências Python do \textit{Hephaestus}.}
\label{tab:deps}
\begin{tabular}{l l p{7cm}}
\hline
\textbf{Pacote} & \textbf{Versão mínima} & \textbf{Função no projeto} \\
\hline
\texttt{sympy}         & 1.12   & Álgebra simbólica, derivação de $M$, $C$, $G$ e \texttt{lambdify} \\
\texttt{numpy}         & 1.24   & Álgebra linear numérica, avaliação das funções compiladas \\
\texttt{matplotlib}    & 3.7    & Visualização dos gráficos de erro, torque e trajetória \\
\texttt{customtkinter} & 5.2    & Interface gráfica (\textit{GUI}) com tema moderno \\
\texttt{scipy}         & 1.11   & Funções auxiliares (interpolação, transformadas) \\
\hline
\end{tabular}
\end{table}

Todas as dependências listadas são pacotes de código aberto disponíveis no repositório PyPI.

% ------------------------------------------------------------------------------
\section*{C.3 Instalação}
% ------------------------------------------------------------------------------

Recomenda-se a criação de um ambiente virtual Python isolado para evitar conflitos entre versões de pacotes:

\begin{lstlisting}[style=bash, caption={Criação de ambiente virtual e instalação das dependências.}, label={lst:install}]
# Clonar o repositorio (ou extrair o arquivo comprimido)
git clone https://github.com/seu_usuario/hephaestus.git
cd hephaestus

# Criar e ativar ambiente virtual (Windows)
python -m venv .venv
.venv\Scripts\activate

# Instalar dependencias
pip install sympy numpy matplotlib customtkinter scipy

# Verificar instalacao
python -c "import sympy, numpy, matplotlib, customtkinter; print('OK')"
\end{lstlisting}

Para sistemas Linux/macOS, substituir \texttt{.venv\textbackslash Scripts\textbackslash activate} por \texttt{source .venv/bin/activate}.

% ------------------------------------------------------------------------------
\section*{C.4 Execução}
% ------------------------------------------------------------------------------

O ponto de entrada principal da aplicação é o arquivo \texttt{main.py}. Para iniciar a interface gráfica:

\begin{lstlisting}[style=bash, caption={Execução da interface gráfica do Hephaestus.}, label={lst:run}]
# Com o ambiente virtual ativado
python main.py
\end{lstlisting}

\textbf{Importante (Windows):} o \textit{Hephaestus} utiliza \texttt{ProcessPoolExecutor} com método de \textit{spawn} para o paralelismo de processos. Em sistemas Windows, é obrigatório que o ponto de entrada esteja protegido pelo bloco \texttt{if \_\_name\_\_ == '\_\_main\_\_':}, conforme implementado em \texttt{main.py}:

\begin{lstlisting}[language=Python, caption={Bloco de proteção obrigatório no \texttt{main.py} para Windows.}, label={lst:main_guard}]
import multiprocessing

if __name__ == '__main__':
    multiprocessing.freeze_support()   # necessario para executaveis PyInstaller
    app = App()
    app.mainloop()
\end{lstlisting}

A ausência deste bloco causaria a abertura recursiva de múltiplas janelas ao criar processos filhos no Windows.

% ------------------------------------------------------------------------------
\section*{C.5 Geração do Executável Autônomo (\textit{.exe})}
% ------------------------------------------------------------------------------

Para distribuir o \textit{Hephaestus} como um executável autônomo sem necessidade de instalação de Python, utiliza-se o \textit{PyInstaller}:

\begin{lstlisting}[style=bash, caption={Geração do executável PyInstaller.}, label={lst:pyinstaller}]
pip install pyinstaller

pyinstaller --onefile --windowed \
            --add-data "tooltip_utils.py;." \
            --name "Hephaestus" main.py
\end{lstlisting}

O arquivo resultante \texttt{dist/Hephaestus.exe} é executável sem instalação prévia de Python. Nesse modo, o suporte a \textit{spawn} de processos no Windows é detectado automaticamente pela condição \texttt{sys.frozen} do PyInstaller: quando o executável está \textit{frozen}, o número de \textit{workers} do \texttt{ProcessPoolExecutor} é automaticamente definido como 1 (modo serial), evitando o conflito de \textit{spawn} em módulos congelados. A derivação simbólica e a simulação operam normalmente, apenas sem paralelismo de processos.

% ------------------------------------------------------------------------------
\section*{C.6 Estrutura de Arquivos do Projeto}
% ------------------------------------------------------------------------------

\begin{table}[h!]
\centering
\caption{Estrutura principal de arquivos do repositório \textit{Hephaestus}.}
\label{tab:estrutura}
\begin{tabular}{l p{9.5cm}}
\hline
\textbf{Arquivo} & \textbf{Descrição} \\
\hline
\texttt{main.py}          & Ponto de entrada; interface gráfica (\textit{CustomTkinter}) com as abas Modelagem, Simulação e Análise \\
\texttt{engine.py}        & \textit{Engine} matemática: classes \texttt{RobotMathEngine} e \texttt{RobotMathHydro}; derivação simbólica de $M$, $C$, $G$, $J$ \\
\texttt{simulator.py}     & Simulador numérico: classe \texttt{RobotSimulator}; controladores \texttt{SlidingModeControl}, \texttt{SuperTwistingSMC}, \texttt{Decentralized\_LADRC}; integrador simplético \\
\texttt{tooltip\_utils.py} & Utilitários de interface: tooltips educacionais e blocos de ajuda contextual para os parâmetros físicos \\
\texttt{TCC/}             & Capítulos do documento TCC em \LaTeX{} \\
\hline
\end{tabular}
\end{table}

%
% Pacotes necessários no preâmbulo do documento principal:
%   \usepackage{listings}
%   \usepackage{xcolor}
% ==============================================================================

% Configuração global das listagens de código Python
\lstset{
  language=Python,
  basicstyle=\ttfamily\footnotesize,
  keywordstyle=\color{blue}\bfseries,
  commentstyle=\color{gray}\itshape,
  stringstyle=\color{teal},
  numberstyle=\tiny\color{gray},
  numbers=left,
  stepnumber=1,
  numbersep=8pt,
  backgroundcolor=\color{gray!6},
  frame=single,
  framerule=0.4pt,
  rulecolor=\color{gray!40},
  breaklines=true,
  breakatwhitespace=false,
  showstringspaces=false,
  tabsize=4,
  captionpos=b,
  morekeywords={self, True, False, None, yield, lambda, with, as, assert,
                del, elif, except, finally, from, global, import, in,
                is, not, or, pass, raise, try},
  literate={á}{{\'a}}1 {é}{{\'e}}1 {í}{{\'i}}1 {ó}{{\'o}}1 {ú}{{\'u}}1
           {â}{{\^a}}1 {ê}{{\^e}}1 {î}{{\^i}}1 {ô}{{\^o}}1 {û}{{\^u}}1
           {ã}{{\~a}}1 {õ}{{\~o}}1 {ç}{{\c{c}}}1
           {Á}{{\'A}}1 {É}{{\'E}}1 {Í}{{\'I}}1 {Ó}{{\'O}}1 {Ú}{{\'U}}1
           {Â}{{\^A}}1 {Ê}{{\^E}}1 {Ô}{{\^O}}1 {Ã}{{\~A}}1 {Õ}{{\~O}}1
           {Ç}{{\c{C}}}1,
}

\lstdefinestyle{bash}{
  language=bash,
  keywordstyle=\color{violet}\bfseries,
  commentstyle=\color{gray}\itshape,
  stringstyle=\color{teal},
  backgroundcolor=\color{gray!8},
}

\chapter{Código-fonte da \textit{Engine} Matemática (\texttt{engine.py})}
\label{ap:engine}

Este apêndice apresenta o código-fonte completo dos componentes fundamentais da \textit{engine} matemática do \textit{Hephaestus}, implementados no módulo \texttt{engine.py}. As listagens cobrem as funções de nível de módulo utilizadas para paralelismo, a classe \texttt{RobotMathEngine} (ambiente seco) e a subclasse \texttt{RobotMathHydro} (ambiente subaquático). O código é apresentado na íntegra para viabilizar reprodutibilidade e auditoria matemática das equações derivadas.

% ------------------------------------------------------------------------------
\section*{A.1 Funções de Nível de Módulo para Paralelismo}
% ------------------------------------------------------------------------------

As duas funções a seguir são definidas no escopo do módulo para permitir sua serialização via \texttt{pickle}, requisito do \texttt{ProcessPoolExecutor} no Python. A função \texttt{\_diff\_M\_col} diferencia a matriz de inércia em relação a uma variável de junta $q_k$; a função \texttt{\_calc\_coriolis\_row} monta a $i$-ésima linha do vetor de Coriolis pela fórmula dos símbolos de Christoffel.

\begin{lstlisting}[language=Python, caption={Funções paralelas para derivação de $\partial M/\partial q_k$ e cálculo de $C_i$.}, label={lst:parallel_funcs}]
from concurrent.futures import ProcessPoolExecutor
from itertools import repeat
import os, time
import sympy as sp
from sympy.physics.mechanics import dynamicsymbols

def _diff_M_col(args):
    """Diferencia a matriz M em relação a uma variável simbólica qk.
    Executado em processo separado para paralelizar o cálculo de dM/dq.
    """
    M_list, n, qk = args
    M = sp.Matrix(n, n, M_list)
    return list(M.diff(qk))


def _calc_coriolis_row(i, q, dq, dM_dq):
    n = len(q)
    termo_linha = sp.Integer(0)
    for j in range(n):
        for k in range(j, n):
            dM_ij_dk = dM_dq[k][i, j]
            dM_ik_dj = dM_dq[j][i, k]
            dM_jk_di = dM_dq[i][j, k]
            c_ijk = sp.Rational(1, 2) * (dM_ij_dk + dM_ik_dj - dM_jk_di)
            if c_ijk != 0:
                termo = c_ijk * dq[j] * dq[k]
                if k != j:
                    termo *= 2
                termo_linha += termo
    # sp.collect() omitido intencionalmente: custo superlinear sem benefício,
    # pois o cse=True no lambdify posterior realiza compactação sistematica.
    return termo_linha
\end{lstlisting}

% ------------------------------------------------------------------------------
\section*{A.2 Classe \texttt{RobotMathEngine} — Ambiente Seco}
% ------------------------------------------------------------------------------

\begin{lstlisting}[language=Python, caption={Construtor e configuração de símbolos de \texttt{RobotMathEngine}.}, label={lst:engine_init}]
class RobotMathEngine:
    def __init__(self, joint_config, link_vectors_mask,
                 logger_callback=None, num_workers=None):
        self.joint_config    = joint_config
        self.link_vectors_mask = [sp.Matrix(v) for v in link_vectors_mask]
        self.log             = logger_callback if logger_callback else print
        self.num_workers     = num_workers

        self.t = sp.symbols('t')
        self.g = sp.symbols('g')

        self.q, self.dq, self.params_list   = [], [], []
        self.frames, self.rotation_matrices  = [], []
        self.com_positions_global            = []
        self.angular_velocities              = []
        self.masses                          = []

        self.params_list.append(self.g)   # g sempre presente como parâmetro
        self.M, self.G_vec, self.C_total, self.Jacobian = None, None, None, None

    def _rot_matrix_local(self, axis, angle):
        c, s = sp.cos(angle), sp.sin(angle)
        if axis == 'x': return sp.Matrix([[1,0,0],[0,c,-s],[0,s,c]])
        if axis == 'y': return sp.Matrix([[c,0,s],[0,1,0],[-s,0,c]])
        if axis == 'z': return sp.Matrix([[c,-s,0],[s,c,0],[0,0,1]])
        return sp.eye(3)

    def _get_axis_vector(self, axis):
        if axis == 'x': return sp.Matrix([1, 0, 0])
        if axis == 'y': return sp.Matrix([0, 1, 0])
        if axis == 'z': return sp.Matrix([0, 0, 1])
        return sp.Matrix([0, 0, 0])
\end{lstlisting}

\begin{lstlisting}[language=Python, caption={Cinemática simbólica direta: transformações homogêneas e posições de CM.}, label={lst:step1}]
    def step_1_kinematics(self):
        self.log("1. (AR) Calculando Cinemarica (Com CM Arbitrario)...")
        T_acc   = sp.eye(4)
        R_acc   = sp.eye(3)
        omega_acc = sp.Matrix([0, 0, 0])

        for i, (j_type, link_vec) in enumerate(
                zip(self.joint_config, self.link_vectors_mask)):
            q  = dynamicsymbols(f'q{i+1}')
            dq = dynamicsymbols(f'q{i+1}', 1)
            self.q.append(q);  self.dq.append(dq)

            m  = sp.symbols(f'm{i+1}')
            L  = sp.symbols(f'L{i+1}')
            cx, cy, cz = sp.symbols(f'cx{i+1} cy{i+1} cz{i+1}')

            self.masses.append(m)
            self.params_list.extend([m, L, cx, cy, cz])

            type_char = j_type[0]
            axis_char = j_type[1].lower()

            if type_char == 'R':            # Junta de revolucao
                axis_vec_global = R_acc * self._get_axis_vector(axis_char)
                omega_new = omega_acc + axis_vec_global * dq
                R_j = self._rot_matrix_local(axis_char, q)
                P_j = sp.Matrix([0, 0, 0])
            elif type_char == 'D':          # Junta prismatica
                omega_new = omega_acc
                R_j = sp.eye(3)
                P_j = sp.Matrix([0, 0, 0])
                if axis_char == 'x': P_j[0] = q
                if axis_char == 'y': P_j[1] = q
                if axis_char == 'z': P_j[2] = q

            R_acc = R_acc * R_j
            self.rotation_matrices.append(R_acc)
            self.angular_velocities.append(omega_new)
            omega_acc = omega_new

            T_joint = sp.eye(4)
            T_joint[0:3, 0:3] = R_j
            T_joint[0:3, 3]   = P_j
            T_at_joint_start  = T_acc * T_joint

            P_link_vec        = link_vec * L
            T_link            = sp.eye(4)
            T_link[0:3, 3]    = P_link_vec
            T_acc             = T_at_joint_start * T_link
            self.frames.append(T_acc)

            # CM Global: transforma coordenadas locais pelo frame ate a junta
            v_cm_local   = sp.Matrix([cx, cy, cz, 1])
            p_cm_global  = T_at_joint_start * v_cm_local
            self.com_positions_global.append(p_cm_global[0:3, 0])
\end{lstlisting}

\begin{lstlisting}[language=Python, caption={Derivação simbólica de $M(q)$, $G(q)$ e $J(q)$ pelo método do Jacobiano da inércia.}, label={lst:step2}]
    def step_2_jacobian_M_G(self):
        self.log("2. (AR) Dinamica M e G (Steiner + Tensores)...")
        _t0 = time.perf_counter()
        n   = len(self.q)
        self.M = sp.zeros(n, n)
        V_tot  = 0

        for i in range(n):
            m       = self.masses[i]
            pos_cm  = self.com_positions_global[i]
            R_global = self.rotation_matrices[i]

            J_v = pos_cm.jacobian(self.q)         # Jacobiano linear do CM
            J_w = self.angular_velocities[i].jacobian(self.dq)  # Jacobiano angular

            Ixx, Iyy, Izz = sp.symbols(f'Ixx{i+1} Iyy{i+1} Izz{i+1}')
            self.params_list.extend([Ixx, Iyy, Izz])

            # Transformacao de Steiner: I_global = R * I_local * R^T
            I_local  = sp.Matrix([[Ixx, 0, 0], [0, Iyy, 0], [0, 0, Izz]])
            I_global = R_global * I_local * R_global.T

            # M += m*Jv^T*Jv + Jw^T * I_global * Jw  (metodo do Jacobiano)
            self.M  += m * J_v.T * J_v + J_w.T * I_global * J_w

            V_tot   += m * self.g * pos_cm[2]   # Energia potencial gravitacional
            self.log(f"  M/G: elo {i+1}/{n} concluido")

        self.G_vec   = sp.Matrix([V_tot]).jacobian(self.q).T
        J_lin        = self.frames[-1][0:3, 3].jacobian(self.q)
        J_ang        = self.angular_velocities[-1].jacobian(self.dq)
        self.Jacobian = J_lin.col_join(J_ang)    # Jacobiano 6xN completo
        self.log(f"  M/G/J: concluido em {time.perf_counter()-_t0:.1f}s")
\end{lstlisting}

\begin{lstlisting}[language=Python, caption={Cálculo do vetor de Coriolis pelos símbolos de Christoffel com paralelismo opcional.}, label={lst:step3}]
    def step_3_coriolis_combined(self):
        self.log("3. Calculando Coriolis...")
        n          = len(self.q)
        self.C_total = sp.zeros(n, 1)
        cpu_total  = os.cpu_count() or 1
        workers    = self.num_workers if self.num_workers is not None else cpu_total
        use_parallel = workers is not None and workers > 1

        # --- Passo 3a: dM/dq  (paralelizavel por coluna) ---
        self.log("  dM/dq: diferenciando M" +
                 (" em paralelo" if use_parallel else "") + "...")
        _t0 = time.perf_counter()
        M_list = list(self.M)
        tasks  = [(M_list, n, qk) for qk in self.q]

        if use_parallel:
            with ProcessPoolExecutor(max_workers=workers) as ex:
                dM_dq_flat = list(ex.map(_diff_M_col, tasks))
        else:
            dM_dq_flat = [_diff_M_col(t) for t in tasks]

        dM_dq = [sp.Matrix(n, n, flat) for flat in dM_dq_flat]
        self.log(f"  dM/dq: concluido em {time.perf_counter()-_t0:.1f}s")

        # --- Passo 3b: Simbolos de Christoffel  C_i(q,dq) ---
        if use_parallel:
            with ProcessPoolExecutor(max_workers=workers) as ex:
                for i, row in enumerate(
                    ex.map(_calc_coriolis_row,
                           range(n), repeat(self.q),
                           repeat(self.dq), repeat(dM_dq)),
                    start=1):
                    self.C_total[i - 1] = row
                    self.log(f"  Coriolis: linha {i}/{n} concluida")
        else:
            for i in range(n):
                self.C_total[i] = _calc_coriolis_row(i, self.q, self.dq, dM_dq)
                self.log(f"  Coriolis: linha {i+1}/{n} concluida")
\end{lstlisting}

\begin{lstlisting}[language=Python, caption={Substituição de símbolos dinâmicos e exportação do dicionário de modelo.}, label={lst:step4}]
    def step_4_prepare_export(self):
        self.log("4. Otimizando equacoes...")
        mapa_subs = {}
        for i in range(len(self.q)):
            mapa_subs[self.q[i]]  = sp.Symbol(f'q{i+1}')
            mapa_subs[self.dq[i]] = sp.Symbol(f'dq{i+1}')

        dq_vec      = sp.Matrix(self.dq)
        V_cartesian = self.Jacobian * dq_vec if self.Jacobian is not None else None

        return {
            "M":      self.M,
            "G":      self.G_vec,
            "C":      self.C_total,
            "J":      self.Jacobian,
            "FK_Pos": self.frames[-1],
            "FK_Vel": V_cartesian,
            "Subs":   mapa_subs,
            "Mode":   "Air"
        }
\end{lstlisting}

% ------------------------------------------------------------------------------
\section*{A.3 Subclasse \texttt{RobotMathHydro} — Extensão Hidrodinâmica}
% ------------------------------------------------------------------------------

\begin{lstlisting}[language=Python, caption={Extensão hidrodinâmica: massa adicionada, empuxo e potencial aparente.}, label={lst:hydro}]
class RobotMathHydro(RobotMathEngine):
    def __init__(self, joint_config, link_vectors_mask,
                 logger_callback=None, num_workers=None):
        super().__init__(joint_config, link_vectors_mask,
                         logger_callback, num_workers)
        self.rho     = sp.symbols('rho')   # densidade do fluido
        self.volumes = []

    def step_1_kinematics_hydro(self):
        super().step_1_kinematics()
        self.log("   -> Adicionando parametros hidrodinamicos...")
        for i in range(len(self.q)):
            vol = sp.symbols(f'vol{i+1}')
            self.volumes.append(vol)
            self.params_list.append(vol)   # vol_i adicionado a assinatura

    def step_2_jacobian_M_G(self):
        self.log("2. (AGUA) Dinamica: M (Inercia + Added Mass) e G (Peso - Empuxo)...")
        _t0  = time.perf_counter()
        n    = len(self.q)
        self.M = sp.zeros(n, n)
        V_tot  = 0

        for i in range(n):
            m       = self.masses[i]
            vol     = self.volumes[i]
            pos_cm  = self.com_positions_global[i]
            R_global = self.rotation_matrices[i]

            J_v = pos_cm.jacobian(self.q)
            J_w = self.angular_velocities[i].jacobian(self.dq)

            # Tensor de inercia rigido (referencial global via Steiner)
            Ixx, Iyy, Izz = sp.symbols(f'Ixx{i+1} Iyy{i+1} Izz{i+1}')
            self.params_list.extend([Ixx, Iyy, Izz])
            I_local_RB  = sp.Matrix([[Ixx,0,0],[0,Iyy,0],[0,0,Izz]])
            I_global_RB = R_global * I_local_RB * R_global.T

            # Tensor de massa adicionada (diagonal local, 6 componentes por elo)
            ma_u, ma_v, ma_w = sp.symbols(f'ma_u{i+1} ma_v{i+1} ma_w{i+1}')
            ma_p, ma_q, ma_r = sp.symbols(f'ma_p{i+1} ma_q{i+1} ma_r{i+1}')
            self.params_list.extend([ma_u, ma_v, ma_w, ma_p, ma_q, ma_r])

            MA_lin_local = sp.Matrix([[ma_u,0,0],[0,ma_v,0],[0,0,ma_w]])
            MA_rot_local = sp.Matrix([[ma_p,0,0],[0,ma_q,0],[0,0,ma_r]])

            # Rotacao para referencial global: M_A^(g) = R * M_A^(l) * R^T
            MA_lin_global = R_global * MA_lin_local * R_global.T
            MA_rot_global = R_global * MA_rot_local * R_global.T

            # M_hydro += Jv^T (m*I + M_A,lin)^(g) Jv + Jw^T (I_RB + M_A,rot)^(g) Jw
            self.M += J_v.T * (m * sp.eye(3) + MA_lin_global) * J_v
            self.M += J_w.T * (I_global_RB + MA_rot_global) * J_w

            # Potencial aparente: V_i = (m - rho*vol) * g * z_cm  (Peso - Empuxo)
            peso_aparente = (m - self.rho * vol) * self.g
            V_tot += peso_aparente * pos_cm[2]
            self.log(f"  M/G: elo {i+1}/{n} concluido")

        self.G_vec   = sp.Matrix([V_tot]).jacobian(self.q).T
        J_lin        = self.frames[-1][0:3, 3].jacobian(self.q)
        J_ang        = self.angular_velocities[-1].jacobian(self.dq)
        self.Jacobian = J_lin.col_join(J_ang)
        self.log(f"  M/G/J: concluido em {time.perf_counter()-_t0:.1f}s")

    def step_4_prepare_export(self):
        data = super().step_4_prepare_export()
        data["Mode"] = "Hydro"
        return data
\end{lstlisting}

% ==============================================================================
\chapter{Código-fonte dos Controladores e do Simulador (\texttt{simulator.py})}
\label{ap:simulador}

Este apêndice apresenta as classes de controle e a estrutura principal da classe \texttt{RobotSimulator} implementadas no módulo \texttt{simulator.py}. As listagens incluem os controladores CT-SMC, Super-Twisting (STA) e LADRC descentralizado, bem como o construtor do simulador, que realiza a compilação das expressões simbólicas via \texttt{lambdify} com CSE, e a infraestrutura de avaliação paralela numérica.

% ------------------------------------------------------------------------------
\section*{B.1 Funções Auxiliares de Orientação}
% ------------------------------------------------------------------------------

\begin{lstlisting}[language=Python, caption={Erro de orientação pelo operador de anel (skew-symmetric) e SLERP em $SO(3)$.}, label={lst:orient_helpers}]
import numpy as np

def _rotation_error_vec(R_ref, R_curr):
    """Erro de orientacao 3D pela formulacao de anel (skew-symmetric).

    Retorna e_R = vee(0.5*(R_ref @ R_curr^T - R_curr @ R_ref^T)),
    proporcional ao eixo-angulo para erros pequenos.
    """
    skew = 0.5 * (R_ref @ R_curr.T - R_curr @ R_ref.T)
    return np.array([skew[2, 1], skew[0, 2], skew[1, 0]])


def _slerp_rotation(R0, Rf, s, sd):
    """SLERP entre R0 e Rf na fracao s em [0, 1].

    Parâmetros
    ----------
    R0, Rf : matrizes de rotacao 3x3
    s      : fracao de interpolacao atual  (0 = R0, 1 = Rf)
    sd     : derivada temporal de s (ds/dt)

    Retorna
    -------
    R_ref     : rotacao interpolada 3x3
    omega_ref : velocidade angular no referencial global (rad/s)
    """
    dR        = R0.T @ Rf
    trace_val = np.clip((np.trace(dR) - 1.0) / 2.0, -1.0, 1.0)
    theta_total = np.arccos(trace_val)
    if abs(theta_total) < 1e-9:
        return R0.copy(), np.zeros(3)
    sin_t      = np.sin(theta_total)
    skew       = (dR - dR.T) / (2.0 * sin_t)
    axis_local = np.array([skew[2, 1], skew[0, 2], skew[1, 0]])
    theta_s    = s * theta_total
    K = np.array([[ 0.0,           -axis_local[2],  axis_local[1]],
                  [ axis_local[2],  0.0,            -axis_local[0]],
                  [-axis_local[1],  axis_local[0],   0.0          ]])
    dR_s  = np.eye(3) + np.sin(theta_s) * K + (1.0 - np.cos(theta_s)) * (K @ K)
    R_ref = R0 @ dR_s
    omega_ref = (sd * theta_total) * (R0 @ axis_local)
    return R_ref, omega_ref
\end{lstlisting}

% ------------------------------------------------------------------------------
\section*{B.2 Controlador CT-SMC com Camada Limite}
% ------------------------------------------------------------------------------

\begin{lstlisting}[language=Python, caption={Controle por Torque Computado com Modo Deslizante (CT-SMC).}, label={lst:smc}]
class SlidingModeControl:
    """Controlador CT-SMC: Computed Torque + Sliding Mode com camada limite.

    Lei de controle:
        q_ddot_f = alpha*q_ddot_d + (1-alpha)*q_ddot_f  (filtro passabaixa)
        v        = q_ddot_f - lambda*e_dot - K*tanh(S/phi)
        tau      = M(q)*v + C(q,dq) + G(q)
    onde S = e_dot + lambda*e  (superficie deslizante).

    Referencia: Slotine & Li, Applied Nonlinear Control, 1991.
    """

    def __init__(self, num_dof, lambda_gain, K, phi, ddq_filter_alpha=1.0):
        self.lambda_gain    = (np.full(num_dof, lambda_gain, dtype=float)
                               if np.isscalar(lambda_gain)
                               else np.array(lambda_gain, dtype=float))
        self.K              = (np.full(num_dof, K, dtype=float)
                               if np.isscalar(K) else np.array(K, dtype=float))
        self.phi            = (np.full(num_dof, phi, dtype=float)
                               if np.isscalar(phi) else np.array(phi, dtype=float))
        self.phi            = np.where(self.phi == 0, 1e-6, self.phi)
        self.alpha          = float(np.clip(ddq_filter_alpha, 1e-6, 1.0))
        self._ddq_d_filt    = None

    def compute_tau(self, q, dq, q_d, dq_d, ddq_d, M, C, G, dt=0.0):
        if self._ddq_d_filt is None:
            self._ddq_d_filt = ddq_d.copy()
        else:
            self._ddq_d_filt = (self.alpha * ddq_d
                                + (1.0 - self.alpha) * self._ddq_d_filt)
        e       = q - q_d
        e_dot   = dq - dq_d
        S       = e_dot + self.lambda_gain * e
        v       = (self._ddq_d_filt - self.lambda_gain * e_dot
                   - self.K * np.tanh(S / self.phi))
        return M @ v + C + G
\end{lstlisting}

% ------------------------------------------------------------------------------
\section*{B.3 Controlador Super-Twisting (STA)}
% ------------------------------------------------------------------------------

\begin{lstlisting}[language=Python, caption={Controlador Super-Twisting de 2ª ordem (STA) com sinal de controle contínuo.}, label={lst:sta}]
class SuperTwistingSMC:
    """Controlador Super-Twisting SMC (STA) baseado em modelo.

    Algoritmo de modo deslizante de 2a ordem:
        q_ddot_f = alpha*q_ddot_d + (1-alpha)*q_ddot_f
        v_sw     = -k1*|S|^(1/2)*sign(S) + z
        z_dot    = -k2*sign(S)
        tau      = M(q)*[q_ddot_f - lambda*e_dot + v_sw] + C(q,dq) + G(q)

    Propriedades: convergencia em tempo finito, sinal de controle continuo.
    Condicoes de ganho (Moreno & Osorio, 2012): k2 > delta, k1 > 2*sqrt(k2*delta).

    Referencias:
    - Levant, Int. J. Control, 58(6), 1993.
    - Moreno & Osorio, IEEE Trans. Autom. Control, 57(4), 2012.
    """

    def __init__(self, num_dof, lambda_gain, k1, k2, ddq_filter_alpha=1.0):
        self.lambda_gain = (np.full(num_dof, lambda_gain, dtype=float)
                            if np.isscalar(lambda_gain)
                            else np.array(lambda_gain, dtype=float))
        self.k1          = (np.full(num_dof, k1, dtype=float)
                            if np.isscalar(k1) else np.array(k1, dtype=float))
        self.k2          = (np.full(num_dof, k2, dtype=float)
                            if np.isscalar(k2) else np.array(k2, dtype=float))
        self.z           = np.zeros(num_dof)
        self.alpha       = float(np.clip(ddq_filter_alpha, 1e-6, 1.0))
        self._ddq_d_filt = None

    def reset(self):
        self.z[:]        = 0.0
        self._ddq_d_filt = None

    def compute_tau(self, q, dq, q_d, dq_d, ddq_d, M, C, G, dt):
        if self._ddq_d_filt is None:
            self._ddq_d_filt = ddq_d.copy()
        else:
            self._ddq_d_filt = (self.alpha * ddq_d
                                + (1.0 - self.alpha) * self._ddq_d_filt)
        e     = q - q_d
        e_dot = dq - dq_d
        S     = e_dot + self.lambda_gain * e

        # Parcela de chaveamento continua do STA
        v_sw = -self.k1 * np.abs(S) ** 0.5 * np.sign(S) + self.z

        v        = self._ddq_d_filt - self.lambda_gain * e_dot + v_sw
        self.z  -= self.k2 * np.sign(S) * dt    # integracao de z (Euler explicito)
        return M @ v + C + G
\end{lstlisting}

% ------------------------------------------------------------------------------
\section*{B.4 LADRC Descentralizado com ESO de Terceira Ordem}
% ------------------------------------------------------------------------------

\begin{lstlisting}[language=Python, caption={LADRC descentralizado: ESO de 3ª ordem e lei de controle por junta.}, label={lst:ladrc}]
class Decentralized_LADRC:
    """LADRC descentralizado para multiplas juntas (2a ordem por junta).

    Para a junta i: ddot_q_i = b0_i * tau_i + f_i  (f_i = disturbio total).
    O ESO de 3a ordem (linear) estima [q_i, dq_i, f_i] em z1, z2, z3.
    Parametrizacao dos ganhos (Gao, 2003): beta1=3*wo, beta2=3*wo^2, beta3=wo^3.
    """

    def __init__(self, num_dof, b0, wo, kp, kd, dt,
                 z_limit=None, tau_limit=None,
                 max_wo_dt=0.1, z3_filter_alpha=1.0):
        self.b0  = (np.full(num_dof, b0, dtype=float)
                    if np.isscalar(b0) else np.array(b0, dtype=float))
        self.wo  = (np.full(num_dof, wo, dtype=float)
                    if np.isscalar(wo) else np.array(wo, dtype=float))
        self.kp  = (np.full(num_dof, kp, dtype=float)
                    if np.isscalar(kp) else np.array(kp, dtype=float))
        self.kd  = (np.full(num_dof, kd, dtype=float)
                    if np.isscalar(kd) else np.array(kd, dtype=float))
        self.dt  = dt
        self.z_limit        = z_limit
        self.tau_limit      = tau_limit
        self.max_wo_dt      = max_wo_dt
        self.z3_filter_alpha = float(np.clip(z3_filter_alpha, 0.0, 1.0))
        self.z3_filtered    = None

        self._refresh_gains()
        self.z1 = np.zeros(num_dof, dtype=float)
        self.z2 = np.zeros(num_dof, dtype=float)
        self.z3 = np.zeros(num_dof, dtype=float)

    def _refresh_gains(self):
        # Limita wo para garantir estabilidade do ESO discreto: wo*dt <= max_wo_dt
        if self.max_wo_dt is not None:
            wo_dt = self.wo * self.dt
            self.wo = np.where(wo_dt > self.max_wo_dt,
                               self.max_wo_dt / self.dt, self.wo)
        self.beta1 = 3.0 *  self.wo
        self.beta2 = 3.0 * (self.wo ** 2)
        self.beta3 =         self.wo ** 3

    def reset_state(self, q, dq, z3=None):
        self.z1 = np.array(q,  dtype=float).copy()
        self.z2 = np.array(dq, dtype=float).copy()
        self.z3 = (np.zeros(self.num_dof, dtype=float) if z3 is None
                   else (np.full(self.num_dof, z3, dtype=float)
                         if np.isscalar(z3) else np.array(z3, dtype=float)))
        self.z3_filtered = self.z3.copy()

    def update_eso(self, q, tau_prev):
        """Avanca o ESO por um passo de integracao (Euler explicito)."""
        error  = self.z1 - q
        z1_dot = self.z2 - self.beta1 * error
        z2_dot = self.z3 - self.beta2 * error + self.b0 * tau_prev
        z3_dot =         - self.beta3 * error

        self.z1 += z1_dot * self.dt
        self.z2 += z2_dot * self.dt
        self.z3 += z3_dot * self.dt

        # Filtro IIR de 1a ordem em z3 (atenua alta frequencia antes da lei)
        a = self.z3_filter_alpha
        self.z3_filtered = a * self.z3 + (1.0 - a) * self.z3_filtered

    def compute_control(self, q_d, dq_d, ddq_d, q_meas=None, dq_meas=None):
        """Calcula o torque de controle LADRC (cancela disturbio + PD + ff)."""
        q_fb  = q_meas  if q_meas  is not None else self.z1
        dq_fb = dq_meas if dq_meas is not None else self.z2

        z3_use = self.z3_filtered if self.z3_filtered is not None else self.z3
        u0     = ddq_d + self.kp * (q_d - q_fb) + self.kd * (dq_d - dq_fb)
        b0_safe = np.where(self.b0 == 0.0, 1e-6, self.b0)
        u = (u0 - z3_use) / b0_safe
        if self.tau_limit is not None:
            u = np.clip(u, -self.tau_limit, self.tau_limit)
        return u
\end{lstlisting}

% ------------------------------------------------------------------------------
\section*{B.5 Compilação Simbólico-Numérica e Executor Persistente}
% ------------------------------------------------------------------------------

\begin{lstlisting}[language=Python, caption={Funções de worker para avaliação paralela de $M$, $C$, $G$ no executor persistente.}, label={lst:workers}]
import sympy as sp

_WORKER_FUNCS = {}   # dicionario de modulo: preenchido por _init_worker


def _init_worker(sym_vars, expr_M, expr_C, expr_G):
    """Inicializador do worker: compila lambdify localmente no processo."""
    global _WORKER_FUNCS
    _WORKER_FUNCS = {
        "M": sp.lambdify(sym_vars, expr_M, modules="numpy"),
        "C": sp.lambdify(sym_vars, expr_C, modules="numpy"),
        "G": sp.lambdify(sym_vars, expr_G, modules="numpy"),
    }


def _eval_worker(task):
    """Avalia uma das funcoes compiladas com os argumentos numericos fornecidos."""
    func_name, args = task
    return _WORKER_FUNCS[func_name](*args)
\end{lstlisting}

\begin{lstlisting}[language=Python, caption={Construtor de \texttt{RobotSimulator}: compilação \texttt{lambdify} + CSE e inicialização do executor.}, label={lst:sim_init}]
from concurrent.futures import ProcessPoolExecutor
import time, numpy as np, sympy as sp

class RobotSimulator:
    def __init__(self, robot_math_instance, mode="Air"):
        self.bot       = robot_math_instance
        self.mode      = mode
        self.num_dof   = len(self.bot.q)
        self.q_home    = np.zeros(self.num_dof)
        self._executor = None      # executor persistente (inicializado apos compilacao)
        self._executor_workers = 0

        print(f"[{mode}] Compilando equacoes com lambdify+CSE...")
        _t0 = time.perf_counter()

        # Identificacao de juntas rotacionais (para wrap em [-pi, pi])
        self.is_rotational = np.array(
            [j.startswith('R') for j in self.bot.joint_config], dtype=bool)

        # Assinatura completa: [q1..qN, dq1..dqN, params...]
        self.sym_vars = self.bot.q + self.bot.dq + self.bot.params_list
        if hasattr(self.bot, 'rho'):
            self.sym_vars.append(self.bot.rho)

        # Compilacao separada por grandeza (M, C, G, J, FK)
        # cse=True: CSE antes da geracao do codigo -> reducao de 60-90% no bytecode
        self.expr_M = self.bot.M
        self.expr_C = self.bot.C_total
        self.expr_G = self.bot.G_vec

        self.func_M = sp.lambdify(self.sym_vars, self.expr_M,
                                  modules="numpy", cse=True)
        self.func_C = sp.lambdify(self.sym_vars, self.expr_C,
                                  modules="numpy", cse=True)
        self.func_G = sp.lambdify(self.sym_vars, self.expr_G,
                                  modules="numpy", cse=True)
        self.func_J = sp.lambdify(self.sym_vars, self.bot.Jacobian,
                                  modules="numpy", cse=True)

        # Funcoes FK por elo para visualizacao animada
        self.funcs_fk_all_links = [
            sp.lambdify(self.sym_vars,
                        frame[0:3, 3].subs(self.bot.step_4_prepare_export()["Subs"]),
                        modules="numpy", cse=True)
            for frame in self.bot.frames
        ]
        print(f"Compilacao concluida em {time.perf_counter()-_t0:.1f}s")

    def warm_up_workers(self, num_workers):
        """Inicializa o executor persistente e aquece os workers."""
        if (self._executor is not None
                and self._executor_workers == num_workers):
            return   # ja inicializado com mesmo numero de workers
        if self._executor is not None:
            self._executor.shutdown(wait=False)
        self._executor = ProcessPoolExecutor(
            max_workers=num_workers,
            initializer=_init_worker,
            initargs=(self.sym_vars, self.expr_M, self.expr_C, self.expr_G))
        self._executor_workers = num_workers
        # Warmup: submete tarefas dummy para pre-compilar os workers
        dummy_args = tuple(np.zeros(len(self.sym_vars)))
        futs = [self._executor.submit(_eval_worker, (k, dummy_args))
                for k in ("M", "C", "G")]
        for f in futs:
            try: f.result()
            except Exception: pass

    def close(self):
        """Encerra o executor persistente de forma limpa."""
        if self._executor is not None:
            self._executor.shutdown(wait=False)
            self._executor = None
\end{lstlisting}

% ==============================================================================
\chapter{Requisitos de Software, Instalação e Execução do \textit{Hephaestus}}
\label{ap:instalacao}

Este apêndice descreve os requisitos de software, o procedimento de instalação e as instruções de execução do ambiente \textit{Hephaestus} a partir do código-fonte.

% ------------------------------------------------------------------------------
\section*{C.1 Requisitos de Sistema}
% ------------------------------------------------------------------------------

O \textit{Hephaestus} foi desenvolvido e validado no seguinte ambiente:

\begin{itemize}
    \item \textbf{Sistema Operacional:} Windows 10/11 (64 bits). O código é multiplataforma e compatível com Linux e macOS com as mesmas dependências; no entanto, a interface gráfica com \textit{CustomTkinter} pode apresentar diferenças visuais menores entre plataformas.
    \item \textbf{Python:} versão 3.10 ou superior (recomendado: 3.11). O uso de \texttt{ProcessPoolExecutor} com método de \textit{spawn} requer Python $\geq$ 3.8; a tipagem e as dependências utilizadas requerem Python $\geq$ 3.10.
    \item \textbf{Hardware recomendado:} processador com pelo menos 4 núcleos lógicos (para aproveitar o paralelismo de derivação); 8 GB de RAM (para modelos com $N \geq 6$ GDL, cujas matrizes simbólicas podem ocupar centenas de MB antes da compilação).
\end{itemize}

% ------------------------------------------------------------------------------
\section*{C.2 Dependências Python}
% ------------------------------------------------------------------------------

A Tabela~\ref{tab:deps} lista as dependências diretas do projeto com as versões mínimas recomendadas.

\begin{table}[h!]
\centering
\caption{Dependências Python do \textit{Hephaestus}.}
\label{tab:deps}
\begin{tabular}{l l p{7cm}}
\hline
\textbf{Pacote} & \textbf{Versão mínima} & \textbf{Função no projeto} \\
\hline
\texttt{sympy}         & 1.12   & Álgebra simbólica, derivação de $M$, $C$, $G$ e \texttt{lambdify} \\
\texttt{numpy}         & 1.24   & Álgebra linear numérica, avaliação das funções compiladas \\
\texttt{matplotlib}    & 3.7    & Visualização dos gráficos de erro, torque e trajetória \\
\texttt{customtkinter} & 5.2    & Interface gráfica (\textit{GUI}) com tema moderno \\
\texttt{scipy}         & 1.11   & Funções auxiliares (interpolação, transformadas) \\
\hline
\end{tabular}
\end{table}

Todas as dependências listadas são pacotes de código aberto disponíveis no repositório PyPI.

% ------------------------------------------------------------------------------
\section*{C.3 Instalação}
% ------------------------------------------------------------------------------

Recomenda-se a criação de um ambiente virtual Python isolado para evitar conflitos entre versões de pacotes:

\begin{lstlisting}[style=bash, caption={Criação de ambiente virtual e instalação das dependências.}, label={lst:install}]
# Clonar o repositorio (ou extrair o arquivo comprimido)
git clone https://github.com/seu_usuario/hephaestus.git
cd hephaestus

# Criar e ativar ambiente virtual (Windows)
python -m venv .venv
.venv\Scripts\activate

# Instalar dependencias
pip install sympy numpy matplotlib customtkinter scipy

# Verificar instalacao
python -c "import sympy, numpy, matplotlib, customtkinter; print('OK')"
\end{lstlisting}

Para sistemas Linux/macOS, substituir \texttt{.venv\textbackslash Scripts\textbackslash activate} por \texttt{source .venv/bin/activate}.

% ------------------------------------------------------------------------------
\section*{C.4 Execução}
% ------------------------------------------------------------------------------

O ponto de entrada principal da aplicação é o arquivo \texttt{main.py}. Para iniciar a interface gráfica:

\begin{lstlisting}[style=bash, caption={Execução da interface gráfica do Hephaestus.}, label={lst:run}]
# Com o ambiente virtual ativado
python main.py
\end{lstlisting}

\textbf{Importante (Windows):} o \textit{Hephaestus} utiliza \texttt{ProcessPoolExecutor} com método de \textit{spawn} para o paralelismo de processos. Em sistemas Windows, é obrigatório que o ponto de entrada esteja protegido pelo bloco \texttt{if \_\_name\_\_ == '\_\_main\_\_':}, conforme implementado em \texttt{main.py}:

\begin{lstlisting}[language=Python, caption={Bloco de proteção obrigatório no \texttt{main.py} para Windows.}, label={lst:main_guard}]
import multiprocessing

if __name__ == '__main__':
    multiprocessing.freeze_support()   # necessario para executaveis PyInstaller
    app = App()
    app.mainloop()
\end{lstlisting}

A ausência deste bloco causaria a abertura recursiva de múltiplas janelas ao criar processos filhos no Windows.

% ------------------------------------------------------------------------------
\section*{C.5 Geração do Executável Autônomo (\textit{.exe})}
% ------------------------------------------------------------------------------

Para distribuir o \textit{Hephaestus} como um executável autônomo sem necessidade de instalação de Python, utiliza-se o \textit{PyInstaller}:

\begin{lstlisting}[style=bash, caption={Geração do executável PyInstaller.}, label={lst:pyinstaller}]
pip install pyinstaller

pyinstaller --onefile --windowed \
            --add-data "tooltip_utils.py;." \
            --name "Hephaestus" main.py
\end{lstlisting}

O arquivo resultante \texttt{dist/Hephaestus.exe} é executável sem instalação prévia de Python. Nesse modo, o suporte a \textit{spawn} de processos no Windows é detectado automaticamente pela condição \texttt{sys.frozen} do PyInstaller: quando o executável está \textit{frozen}, o número de \textit{workers} do \texttt{ProcessPoolExecutor} é automaticamente definido como 1 (modo serial), evitando o conflito de \textit{spawn} em módulos congelados. A derivação simbólica e a simulação operam normalmente, apenas sem paralelismo de processos.

% ------------------------------------------------------------------------------
\section*{C.6 Estrutura de Arquivos do Projeto}
% ------------------------------------------------------------------------------

\begin{table}[h!]
\centering
\caption{Estrutura principal de arquivos do repositório \textit{Hephaestus}.}
\label{tab:estrutura}
\begin{tabular}{l p{9.5cm}}
\hline
\textbf{Arquivo} & \textbf{Descrição} \\
\hline
\texttt{main.py}          & Ponto de entrada; interface gráfica (\textit{CustomTkinter}) com as abas Modelagem, Simulação e Análise \\
\texttt{engine.py}        & \textit{Engine} matemática: classes \texttt{RobotMathEngine} e \texttt{RobotMathHydro}; derivação simbólica de $M$, $C$, $G$, $J$ \\
\texttt{simulator.py}     & Simulador numérico: classe \texttt{RobotSimulator}; controladores \texttt{SlidingModeControl}, \texttt{SuperTwistingSMC}, \texttt{Decentralized\_LADRC}; integrador simplético \\
\texttt{tooltip\_utils.py} & Utilitários de interface: tooltips educacionais e blocos de ajuda contextual para os parâmetros físicos \\
\texttt{TCC/}             & Capítulos do documento TCC em \LaTeX{} \\
\hline
\end{tabular}
\end{table}

%
% Pacotes necessários no preâmbulo do documento principal:
%   \usepackage{listings}
%   \usepackage{xcolor}
% ==============================================================================

% Configuração global das listagens de código Python
\lstset{
  language=Python,
  basicstyle=\ttfamily\footnotesize,
  keywordstyle=\color{blue}\bfseries,
  commentstyle=\color{gray}\itshape,
  stringstyle=\color{teal},
  numberstyle=\tiny\color{gray},
  numbers=left,
  stepnumber=1,
  numbersep=8pt,
  backgroundcolor=\color{gray!6},
  frame=single,
  framerule=0.4pt,
  rulecolor=\color{gray!40},
  breaklines=true,
  breakatwhitespace=false,
  showstringspaces=false,
  tabsize=4,
  captionpos=b,
  morekeywords={self, True, False, None, yield, lambda, with, as, assert,
                del, elif, except, finally, from, global, import, in,
                is, not, or, pass, raise, try},
  literate={á}{{\'a}}1 {é}{{\'e}}1 {í}{{\'i}}1 {ó}{{\'o}}1 {ú}{{\'u}}1
           {â}{{\^a}}1 {ê}{{\^e}}1 {î}{{\^i}}1 {ô}{{\^o}}1 {û}{{\^u}}1
           {ã}{{\~a}}1 {õ}{{\~o}}1 {ç}{{\c{c}}}1
           {Á}{{\'A}}1 {É}{{\'E}}1 {Í}{{\'I}}1 {Ó}{{\'O}}1 {Ú}{{\'U}}1
           {Â}{{\^A}}1 {Ê}{{\^E}}1 {Ô}{{\^O}}1 {Ã}{{\~A}}1 {Õ}{{\~O}}1
           {Ç}{{\c{C}}}1,
}

\lstdefinestyle{bash}{
  language=bash,
  keywordstyle=\color{violet}\bfseries,
  commentstyle=\color{gray}\itshape,
  stringstyle=\color{teal},
  backgroundcolor=\color{gray!8},
}

\chapter{Código-fonte da \textit{Engine} Matemática (\texttt{engine.py})}
\label{ap:engine}

Este apêndice apresenta o código-fonte completo dos componentes fundamentais da \textit{engine} matemática do \textit{Hephaestus}, implementados no módulo \texttt{engine.py}. As listagens cobrem as funções de nível de módulo utilizadas para paralelismo, a classe \texttt{RobotMathEngine} (ambiente seco) e a subclasse \texttt{RobotMathHydro} (ambiente subaquático). O código é apresentado na íntegra para viabilizar reprodutibilidade e auditoria matemática das equações derivadas.

% ------------------------------------------------------------------------------
\section*{A.1 Funções de Nível de Módulo para Paralelismo}
% ------------------------------------------------------------------------------

As duas funções a seguir são definidas no escopo do módulo para permitir sua serialização via \texttt{pickle}, requisito do \texttt{ProcessPoolExecutor} no Python. A função \texttt{\_diff\_M\_col} diferencia a matriz de inércia em relação a uma variável de junta $q_k$; a função \texttt{\_calc\_coriolis\_row} monta a $i$-ésima linha do vetor de Coriolis pela fórmula dos símbolos de Christoffel.

\begin{lstlisting}[language=Python, caption={Funções paralelas para derivação de $\partial M/\partial q_k$ e cálculo de $C_i$.}, label={lst:parallel_funcs}]
from concurrent.futures import ProcessPoolExecutor
from itertools import repeat
import os, time
import sympy as sp
from sympy.physics.mechanics import dynamicsymbols

def _diff_M_col(args):
    """Diferencia a matriz M em relação a uma variável simbólica qk.
    Executado em processo separado para paralelizar o cálculo de dM/dq.
    """
    M_list, n, qk = args
    M = sp.Matrix(n, n, M_list)
    return list(M.diff(qk))


def _calc_coriolis_row(i, q, dq, dM_dq):
    n = len(q)
    termo_linha = sp.Integer(0)
    for j in range(n):
        for k in range(j, n):
            dM_ij_dk = dM_dq[k][i, j]
            dM_ik_dj = dM_dq[j][i, k]
            dM_jk_di = dM_dq[i][j, k]
            c_ijk = sp.Rational(1, 2) * (dM_ij_dk + dM_ik_dj - dM_jk_di)
            if c_ijk != 0:
                termo = c_ijk * dq[j] * dq[k]
                if k != j:
                    termo *= 2
                termo_linha += termo
    # sp.collect() omitido intencionalmente: custo superlinear sem benefício,
    # pois o cse=True no lambdify posterior realiza compactação sistematica.
    return termo_linha
\end{lstlisting}

% ------------------------------------------------------------------------------
\section*{A.2 Classe \texttt{RobotMathEngine} — Ambiente Seco}
% ------------------------------------------------------------------------------

\begin{lstlisting}[language=Python, caption={Construtor e configuração de símbolos de \texttt{RobotMathEngine}.}, label={lst:engine_init}]
class RobotMathEngine:
    def __init__(self, joint_config, link_vectors_mask,
                 logger_callback=None, num_workers=None):
        self.joint_config    = joint_config
        self.link_vectors_mask = [sp.Matrix(v) for v in link_vectors_mask]
        self.log             = logger_callback if logger_callback else print
        self.num_workers     = num_workers

        self.t = sp.symbols('t')
        self.g = sp.symbols('g')

        self.q, self.dq, self.params_list   = [], [], []
        self.frames, self.rotation_matrices  = [], []
        self.com_positions_global            = []
        self.angular_velocities              = []
        self.masses                          = []

        self.params_list.append(self.g)   # g sempre presente como parâmetro
        self.M, self.G_vec, self.C_total, self.Jacobian = None, None, None, None

    def _rot_matrix_local(self, axis, angle):
        c, s = sp.cos(angle), sp.sin(angle)
        if axis == 'x': return sp.Matrix([[1,0,0],[0,c,-s],[0,s,c]])
        if axis == 'y': return sp.Matrix([[c,0,s],[0,1,0],[-s,0,c]])
        if axis == 'z': return sp.Matrix([[c,-s,0],[s,c,0],[0,0,1]])
        return sp.eye(3)

    def _get_axis_vector(self, axis):
        if axis == 'x': return sp.Matrix([1, 0, 0])
        if axis == 'y': return sp.Matrix([0, 1, 0])
        if axis == 'z': return sp.Matrix([0, 0, 1])
        return sp.Matrix([0, 0, 0])
\end{lstlisting}

\begin{lstlisting}[language=Python, caption={Cinemática simbólica direta: transformações homogêneas e posições de CM.}, label={lst:step1}]
    def step_1_kinematics(self):
        self.log("1. (AR) Calculando Cinemarica (Com CM Arbitrario)...")
        T_acc   = sp.eye(4)
        R_acc   = sp.eye(3)
        omega_acc = sp.Matrix([0, 0, 0])

        for i, (j_type, link_vec) in enumerate(
                zip(self.joint_config, self.link_vectors_mask)):
            q  = dynamicsymbols(f'q{i+1}')
            dq = dynamicsymbols(f'q{i+1}', 1)
            self.q.append(q);  self.dq.append(dq)

            m  = sp.symbols(f'm{i+1}')
            L  = sp.symbols(f'L{i+1}')
            cx, cy, cz = sp.symbols(f'cx{i+1} cy{i+1} cz{i+1}')

            self.masses.append(m)
            self.params_list.extend([m, L, cx, cy, cz])

            type_char = j_type[0]
            axis_char = j_type[1].lower()

            if type_char == 'R':            # Junta de revolucao
                axis_vec_global = R_acc * self._get_axis_vector(axis_char)
                omega_new = omega_acc + axis_vec_global * dq
                R_j = self._rot_matrix_local(axis_char, q)
                P_j = sp.Matrix([0, 0, 0])
            elif type_char == 'D':          # Junta prismatica
                omega_new = omega_acc
                R_j = sp.eye(3)
                P_j = sp.Matrix([0, 0, 0])
                if axis_char == 'x': P_j[0] = q
                if axis_char == 'y': P_j[1] = q
                if axis_char == 'z': P_j[2] = q

            R_acc = R_acc * R_j
            self.rotation_matrices.append(R_acc)
            self.angular_velocities.append(omega_new)
            omega_acc = omega_new

            T_joint = sp.eye(4)
            T_joint[0:3, 0:3] = R_j
            T_joint[0:3, 3]   = P_j
            T_at_joint_start  = T_acc * T_joint

            P_link_vec        = link_vec * L
            T_link            = sp.eye(4)
            T_link[0:3, 3]    = P_link_vec
            T_acc             = T_at_joint_start * T_link
            self.frames.append(T_acc)

            # CM Global: transforma coordenadas locais pelo frame ate a junta
            v_cm_local   = sp.Matrix([cx, cy, cz, 1])
            p_cm_global  = T_at_joint_start * v_cm_local
            self.com_positions_global.append(p_cm_global[0:3, 0])
\end{lstlisting}

\begin{lstlisting}[language=Python, caption={Derivação simbólica de $M(q)$, $G(q)$ e $J(q)$ pelo método do Jacobiano da inércia.}, label={lst:step2}]
    def step_2_jacobian_M_G(self):
        self.log("2. (AR) Dinamica M e G (Steiner + Tensores)...")
        _t0 = time.perf_counter()
        n   = len(self.q)
        self.M = sp.zeros(n, n)
        V_tot  = 0

        for i in range(n):
            m       = self.masses[i]
            pos_cm  = self.com_positions_global[i]
            R_global = self.rotation_matrices[i]

            J_v = pos_cm.jacobian(self.q)         # Jacobiano linear do CM
            J_w = self.angular_velocities[i].jacobian(self.dq)  # Jacobiano angular

            Ixx, Iyy, Izz = sp.symbols(f'Ixx{i+1} Iyy{i+1} Izz{i+1}')
            self.params_list.extend([Ixx, Iyy, Izz])

            # Transformacao de Steiner: I_global = R * I_local * R^T
            I_local  = sp.Matrix([[Ixx, 0, 0], [0, Iyy, 0], [0, 0, Izz]])
            I_global = R_global * I_local * R_global.T

            # M += m*Jv^T*Jv + Jw^T * I_global * Jw  (metodo do Jacobiano)
            self.M  += m * J_v.T * J_v + J_w.T * I_global * J_w

            V_tot   += m * self.g * pos_cm[2]   # Energia potencial gravitacional
            self.log(f"  M/G: elo {i+1}/{n} concluido")

        self.G_vec   = sp.Matrix([V_tot]).jacobian(self.q).T
        J_lin        = self.frames[-1][0:3, 3].jacobian(self.q)
        J_ang        = self.angular_velocities[-1].jacobian(self.dq)
        self.Jacobian = J_lin.col_join(J_ang)    # Jacobiano 6xN completo
        self.log(f"  M/G/J: concluido em {time.perf_counter()-_t0:.1f}s")
\end{lstlisting}

\begin{lstlisting}[language=Python, caption={Cálculo do vetor de Coriolis pelos símbolos de Christoffel com paralelismo opcional.}, label={lst:step3}]
    def step_3_coriolis_combined(self):
        self.log("3. Calculando Coriolis...")
        n          = len(self.q)
        self.C_total = sp.zeros(n, 1)
        cpu_total  = os.cpu_count() or 1
        workers    = self.num_workers if self.num_workers is not None else cpu_total
        use_parallel = workers is not None and workers > 1

        # --- Passo 3a: dM/dq  (paralelizavel por coluna) ---
        self.log("  dM/dq: diferenciando M" +
                 (" em paralelo" if use_parallel else "") + "...")
        _t0 = time.perf_counter()
        M_list = list(self.M)
        tasks  = [(M_list, n, qk) for qk in self.q]

        if use_parallel:
            with ProcessPoolExecutor(max_workers=workers) as ex:
                dM_dq_flat = list(ex.map(_diff_M_col, tasks))
        else:
            dM_dq_flat = [_diff_M_col(t) for t in tasks]

        dM_dq = [sp.Matrix(n, n, flat) for flat in dM_dq_flat]
        self.log(f"  dM/dq: concluido em {time.perf_counter()-_t0:.1f}s")

        # --- Passo 3b: Simbolos de Christoffel  C_i(q,dq) ---
        if use_parallel:
            with ProcessPoolExecutor(max_workers=workers) as ex:
                for i, row in enumerate(
                    ex.map(_calc_coriolis_row,
                           range(n), repeat(self.q),
                           repeat(self.dq), repeat(dM_dq)),
                    start=1):
                    self.C_total[i - 1] = row
                    self.log(f"  Coriolis: linha {i}/{n} concluida")
        else:
            for i in range(n):
                self.C_total[i] = _calc_coriolis_row(i, self.q, self.dq, dM_dq)
                self.log(f"  Coriolis: linha {i+1}/{n} concluida")
\end{lstlisting}

\begin{lstlisting}[language=Python, caption={Substituição de símbolos dinâmicos e exportação do dicionário de modelo.}, label={lst:step4}]
    def step_4_prepare_export(self):
        self.log("4. Otimizando equacoes...")
        mapa_subs = {}
        for i in range(len(self.q)):
            mapa_subs[self.q[i]]  = sp.Symbol(f'q{i+1}')
            mapa_subs[self.dq[i]] = sp.Symbol(f'dq{i+1}')

        dq_vec      = sp.Matrix(self.dq)
        V_cartesian = self.Jacobian * dq_vec if self.Jacobian is not None else None

        return {
            "M":      self.M,
            "G":      self.G_vec,
            "C":      self.C_total,
            "J":      self.Jacobian,
            "FK_Pos": self.frames[-1],
            "FK_Vel": V_cartesian,
            "Subs":   mapa_subs,
            "Mode":   "Air"
        }
\end{lstlisting}

% ------------------------------------------------------------------------------
\section*{A.3 Subclasse \texttt{RobotMathHydro} — Extensão Hidrodinâmica}
% ------------------------------------------------------------------------------

\begin{lstlisting}[language=Python, caption={Extensão hidrodinâmica: massa adicionada, empuxo e potencial aparente.}, label={lst:hydro}]
class RobotMathHydro(RobotMathEngine):
    def __init__(self, joint_config, link_vectors_mask,
                 logger_callback=None, num_workers=None):
        super().__init__(joint_config, link_vectors_mask,
                         logger_callback, num_workers)
        self.rho     = sp.symbols('rho')   # densidade do fluido
        self.volumes = []

    def step_1_kinematics_hydro(self):
        super().step_1_kinematics()
        self.log("   -> Adicionando parametros hidrodinamicos...")
        for i in range(len(self.q)):
            vol = sp.symbols(f'vol{i+1}')
            self.volumes.append(vol)
            self.params_list.append(vol)   # vol_i adicionado a assinatura

    def step_2_jacobian_M_G(self):
        self.log("2. (AGUA) Dinamica: M (Inercia + Added Mass) e G (Peso - Empuxo)...")
        _t0  = time.perf_counter()
        n    = len(self.q)
        self.M = sp.zeros(n, n)
        V_tot  = 0

        for i in range(n):
            m       = self.masses[i]
            vol     = self.volumes[i]
            pos_cm  = self.com_positions_global[i]
            R_global = self.rotation_matrices[i]

            J_v = pos_cm.jacobian(self.q)
            J_w = self.angular_velocities[i].jacobian(self.dq)

            # Tensor de inercia rigido (referencial global via Steiner)
            Ixx, Iyy, Izz = sp.symbols(f'Ixx{i+1} Iyy{i+1} Izz{i+1}')
            self.params_list.extend([Ixx, Iyy, Izz])
            I_local_RB  = sp.Matrix([[Ixx,0,0],[0,Iyy,0],[0,0,Izz]])
            I_global_RB = R_global * I_local_RB * R_global.T

            # Tensor de massa adicionada (diagonal local, 6 componentes por elo)
            ma_u, ma_v, ma_w = sp.symbols(f'ma_u{i+1} ma_v{i+1} ma_w{i+1}')
            ma_p, ma_q, ma_r = sp.symbols(f'ma_p{i+1} ma_q{i+1} ma_r{i+1}')
            self.params_list.extend([ma_u, ma_v, ma_w, ma_p, ma_q, ma_r])

            MA_lin_local = sp.Matrix([[ma_u,0,0],[0,ma_v,0],[0,0,ma_w]])
            MA_rot_local = sp.Matrix([[ma_p,0,0],[0,ma_q,0],[0,0,ma_r]])

            # Rotacao para referencial global: M_A^(g) = R * M_A^(l) * R^T
            MA_lin_global = R_global * MA_lin_local * R_global.T
            MA_rot_global = R_global * MA_rot_local * R_global.T

            # M_hydro += Jv^T (m*I + M_A,lin)^(g) Jv + Jw^T (I_RB + M_A,rot)^(g) Jw
            self.M += J_v.T * (m * sp.eye(3) + MA_lin_global) * J_v
            self.M += J_w.T * (I_global_RB + MA_rot_global) * J_w

            # Potencial aparente: V_i = (m - rho*vol) * g * z_cm  (Peso - Empuxo)
            peso_aparente = (m - self.rho * vol) * self.g
            V_tot += peso_aparente * pos_cm[2]
            self.log(f"  M/G: elo {i+1}/{n} concluido")

        self.G_vec   = sp.Matrix([V_tot]).jacobian(self.q).T
        J_lin        = self.frames[-1][0:3, 3].jacobian(self.q)
        J_ang        = self.angular_velocities[-1].jacobian(self.dq)
        self.Jacobian = J_lin.col_join(J_ang)
        self.log(f"  M/G/J: concluido em {time.perf_counter()-_t0:.1f}s")

    def step_4_prepare_export(self):
        data = super().step_4_prepare_export()
        data["Mode"] = "Hydro"
        return data
\end{lstlisting}

% ==============================================================================
\chapter{Código-fonte dos Controladores e do Simulador (\texttt{simulator.py})}
\label{ap:simulador}

Este apêndice apresenta as classes de controle e a estrutura principal da classe \texttt{RobotSimulator} implementadas no módulo \texttt{simulator.py}. As listagens incluem os controladores CT-SMC, Super-Twisting (STA) e LADRC descentralizado, bem como o construtor do simulador, que realiza a compilação das expressões simbólicas via \texttt{lambdify} com CSE, e a infraestrutura de avaliação paralela numérica.

% ------------------------------------------------------------------------------
\section*{B.1 Funções Auxiliares de Orientação}
% ------------------------------------------------------------------------------

\begin{lstlisting}[language=Python, caption={Erro de orientação pelo operador de anel (skew-symmetric) e SLERP em $SO(3)$.}, label={lst:orient_helpers}]
import numpy as np

def _rotation_error_vec(R_ref, R_curr):
    """Erro de orientacao 3D pela formulacao de anel (skew-symmetric).

    Retorna e_R = vee(0.5*(R_ref @ R_curr^T - R_curr @ R_ref^T)),
    proporcional ao eixo-angulo para erros pequenos.
    """
    skew = 0.5 * (R_ref @ R_curr.T - R_curr @ R_ref.T)
    return np.array([skew[2, 1], skew[0, 2], skew[1, 0]])


def _slerp_rotation(R0, Rf, s, sd):
    """SLERP entre R0 e Rf na fracao s em [0, 1].

    Parâmetros
    ----------
    R0, Rf : matrizes de rotacao 3x3
    s      : fracao de interpolacao atual  (0 = R0, 1 = Rf)
    sd     : derivada temporal de s (ds/dt)

    Retorna
    -------
    R_ref     : rotacao interpolada 3x3
    omega_ref : velocidade angular no referencial global (rad/s)
    """
    dR        = R0.T @ Rf
    trace_val = np.clip((np.trace(dR) - 1.0) / 2.0, -1.0, 1.0)
    theta_total = np.arccos(trace_val)
    if abs(theta_total) < 1e-9:
        return R0.copy(), np.zeros(3)
    sin_t      = np.sin(theta_total)
    skew       = (dR - dR.T) / (2.0 * sin_t)
    axis_local = np.array([skew[2, 1], skew[0, 2], skew[1, 0]])
    theta_s    = s * theta_total
    K = np.array([[ 0.0,           -axis_local[2],  axis_local[1]],
                  [ axis_local[2],  0.0,            -axis_local[0]],
                  [-axis_local[1],  axis_local[0],   0.0          ]])
    dR_s  = np.eye(3) + np.sin(theta_s) * K + (1.0 - np.cos(theta_s)) * (K @ K)
    R_ref = R0 @ dR_s
    omega_ref = (sd * theta_total) * (R0 @ axis_local)
    return R_ref, omega_ref
\end{lstlisting}

% ------------------------------------------------------------------------------
\section*{B.2 Controlador CT-SMC com Camada Limite}
% ------------------------------------------------------------------------------

\begin{lstlisting}[language=Python, caption={Controle por Torque Computado com Modo Deslizante (CT-SMC).}, label={lst:smc}]
class SlidingModeControl:
    """Controlador CT-SMC: Computed Torque + Sliding Mode com camada limite.

    Lei de controle:
        q_ddot_f = alpha*q_ddot_d + (1-alpha)*q_ddot_f  (filtro passabaixa)
        v        = q_ddot_f - lambda*e_dot - K*tanh(S/phi)
        tau      = M(q)*v + C(q,dq) + G(q)
    onde S = e_dot + lambda*e  (superficie deslizante).

    Referencia: Slotine & Li, Applied Nonlinear Control, 1991.
    """

    def __init__(self, num_dof, lambda_gain, K, phi, ddq_filter_alpha=1.0):
        self.lambda_gain    = (np.full(num_dof, lambda_gain, dtype=float)
                               if np.isscalar(lambda_gain)
                               else np.array(lambda_gain, dtype=float))
        self.K              = (np.full(num_dof, K, dtype=float)
                               if np.isscalar(K) else np.array(K, dtype=float))
        self.phi            = (np.full(num_dof, phi, dtype=float)
                               if np.isscalar(phi) else np.array(phi, dtype=float))
        self.phi            = np.where(self.phi == 0, 1e-6, self.phi)
        self.alpha          = float(np.clip(ddq_filter_alpha, 1e-6, 1.0))
        self._ddq_d_filt    = None

    def compute_tau(self, q, dq, q_d, dq_d, ddq_d, M, C, G, dt=0.0):
        if self._ddq_d_filt is None:
            self._ddq_d_filt = ddq_d.copy()
        else:
            self._ddq_d_filt = (self.alpha * ddq_d
                                + (1.0 - self.alpha) * self._ddq_d_filt)
        e       = q - q_d
        e_dot   = dq - dq_d
        S       = e_dot + self.lambda_gain * e
        v       = (self._ddq_d_filt - self.lambda_gain * e_dot
                   - self.K * np.tanh(S / self.phi))
        return M @ v + C + G
\end{lstlisting}

% ------------------------------------------------------------------------------
\section*{B.3 Controlador Super-Twisting (STA)}
% ------------------------------------------------------------------------------

\begin{lstlisting}[language=Python, caption={Controlador Super-Twisting de 2ª ordem (STA) com sinal de controle contínuo.}, label={lst:sta}]
class SuperTwistingSMC:
    """Controlador Super-Twisting SMC (STA) baseado em modelo.

    Algoritmo de modo deslizante de 2a ordem:
        q_ddot_f = alpha*q_ddot_d + (1-alpha)*q_ddot_f
        v_sw     = -k1*|S|^(1/2)*sign(S) + z
        z_dot    = -k2*sign(S)
        tau      = M(q)*[q_ddot_f - lambda*e_dot + v_sw] + C(q,dq) + G(q)

    Propriedades: convergencia em tempo finito, sinal de controle continuo.
    Condicoes de ganho (Moreno & Osorio, 2012): k2 > delta, k1 > 2*sqrt(k2*delta).

    Referencias:
    - Levant, Int. J. Control, 58(6), 1993.
    - Moreno & Osorio, IEEE Trans. Autom. Control, 57(4), 2012.
    """

    def __init__(self, num_dof, lambda_gain, k1, k2, ddq_filter_alpha=1.0):
        self.lambda_gain = (np.full(num_dof, lambda_gain, dtype=float)
                            if np.isscalar(lambda_gain)
                            else np.array(lambda_gain, dtype=float))
        self.k1          = (np.full(num_dof, k1, dtype=float)
                            if np.isscalar(k1) else np.array(k1, dtype=float))
        self.k2          = (np.full(num_dof, k2, dtype=float)
                            if np.isscalar(k2) else np.array(k2, dtype=float))
        self.z           = np.zeros(num_dof)
        self.alpha       = float(np.clip(ddq_filter_alpha, 1e-6, 1.0))
        self._ddq_d_filt = None

    def reset(self):
        self.z[:]        = 0.0
        self._ddq_d_filt = None

    def compute_tau(self, q, dq, q_d, dq_d, ddq_d, M, C, G, dt):
        if self._ddq_d_filt is None:
            self._ddq_d_filt = ddq_d.copy()
        else:
            self._ddq_d_filt = (self.alpha * ddq_d
                                + (1.0 - self.alpha) * self._ddq_d_filt)
        e     = q - q_d
        e_dot = dq - dq_d
        S     = e_dot + self.lambda_gain * e

        # Parcela de chaveamento continua do STA
        v_sw = -self.k1 * np.abs(S) ** 0.5 * np.sign(S) + self.z

        v        = self._ddq_d_filt - self.lambda_gain * e_dot + v_sw
        self.z  -= self.k2 * np.sign(S) * dt    # integracao de z (Euler explicito)
        return M @ v + C + G
\end{lstlisting}

% ------------------------------------------------------------------------------
\section*{B.4 LADRC Descentralizado com ESO de Terceira Ordem}
% ------------------------------------------------------------------------------

\begin{lstlisting}[language=Python, caption={LADRC descentralizado: ESO de 3ª ordem e lei de controle por junta.}, label={lst:ladrc}]
class Decentralized_LADRC:
    """LADRC descentralizado para multiplas juntas (2a ordem por junta).

    Para a junta i: ddot_q_i = b0_i * tau_i + f_i  (f_i = disturbio total).
    O ESO de 3a ordem (linear) estima [q_i, dq_i, f_i] em z1, z2, z3.
    Parametrizacao dos ganhos (Gao, 2003): beta1=3*wo, beta2=3*wo^2, beta3=wo^3.
    """

    def __init__(self, num_dof, b0, wo, kp, kd, dt,
                 z_limit=None, tau_limit=None,
                 max_wo_dt=0.1, z3_filter_alpha=1.0):
        self.b0  = (np.full(num_dof, b0, dtype=float)
                    if np.isscalar(b0) else np.array(b0, dtype=float))
        self.wo  = (np.full(num_dof, wo, dtype=float)
                    if np.isscalar(wo) else np.array(wo, dtype=float))
        self.kp  = (np.full(num_dof, kp, dtype=float)
                    if np.isscalar(kp) else np.array(kp, dtype=float))
        self.kd  = (np.full(num_dof, kd, dtype=float)
                    if np.isscalar(kd) else np.array(kd, dtype=float))
        self.dt  = dt
        self.z_limit        = z_limit
        self.tau_limit      = tau_limit
        self.max_wo_dt      = max_wo_dt
        self.z3_filter_alpha = float(np.clip(z3_filter_alpha, 0.0, 1.0))
        self.z3_filtered    = None

        self._refresh_gains()
        self.z1 = np.zeros(num_dof, dtype=float)
        self.z2 = np.zeros(num_dof, dtype=float)
        self.z3 = np.zeros(num_dof, dtype=float)

    def _refresh_gains(self):
        # Limita wo para garantir estabilidade do ESO discreto: wo*dt <= max_wo_dt
        if self.max_wo_dt is not None:
            wo_dt = self.wo * self.dt
            self.wo = np.where(wo_dt > self.max_wo_dt,
                               self.max_wo_dt / self.dt, self.wo)
        self.beta1 = 3.0 *  self.wo
        self.beta2 = 3.0 * (self.wo ** 2)
        self.beta3 =         self.wo ** 3

    def reset_state(self, q, dq, z3=None):
        self.z1 = np.array(q,  dtype=float).copy()
        self.z2 = np.array(dq, dtype=float).copy()
        self.z3 = (np.zeros(self.num_dof, dtype=float) if z3 is None
                   else (np.full(self.num_dof, z3, dtype=float)
                         if np.isscalar(z3) else np.array(z3, dtype=float)))
        self.z3_filtered = self.z3.copy()

    def update_eso(self, q, tau_prev):
        """Avanca o ESO por um passo de integracao (Euler explicito)."""
        error  = self.z1 - q
        z1_dot = self.z2 - self.beta1 * error
        z2_dot = self.z3 - self.beta2 * error + self.b0 * tau_prev
        z3_dot =         - self.beta3 * error

        self.z1 += z1_dot * self.dt
        self.z2 += z2_dot * self.dt
        self.z3 += z3_dot * self.dt

        # Filtro IIR de 1a ordem em z3 (atenua alta frequencia antes da lei)
        a = self.z3_filter_alpha
        self.z3_filtered = a * self.z3 + (1.0 - a) * self.z3_filtered

    def compute_control(self, q_d, dq_d, ddq_d, q_meas=None, dq_meas=None):
        """Calcula o torque de controle LADRC (cancela disturbio + PD + ff)."""
        q_fb  = q_meas  if q_meas  is not None else self.z1
        dq_fb = dq_meas if dq_meas is not None else self.z2

        z3_use = self.z3_filtered if self.z3_filtered is not None else self.z3
        u0     = ddq_d + self.kp * (q_d - q_fb) + self.kd * (dq_d - dq_fb)
        b0_safe = np.where(self.b0 == 0.0, 1e-6, self.b0)
        u = (u0 - z3_use) / b0_safe
        if self.tau_limit is not None:
            u = np.clip(u, -self.tau_limit, self.tau_limit)
        return u
\end{lstlisting}

% ------------------------------------------------------------------------------
\section*{B.5 Compilação Simbólico-Numérica e Executor Persistente}
% ------------------------------------------------------------------------------

\begin{lstlisting}[language=Python, caption={Funções de worker para avaliação paralela de $M$, $C$, $G$ no executor persistente.}, label={lst:workers}]
import sympy as sp

_WORKER_FUNCS = {}   # dicionario de modulo: preenchido por _init_worker


def _init_worker(sym_vars, expr_M, expr_C, expr_G):
    """Inicializador do worker: compila lambdify localmente no processo."""
    global _WORKER_FUNCS
    _WORKER_FUNCS = {
        "M": sp.lambdify(sym_vars, expr_M, modules="numpy"),
        "C": sp.lambdify(sym_vars, expr_C, modules="numpy"),
        "G": sp.lambdify(sym_vars, expr_G, modules="numpy"),
    }


def _eval_worker(task):
    """Avalia uma das funcoes compiladas com os argumentos numericos fornecidos."""
    func_name, args = task
    return _WORKER_FUNCS[func_name](*args)
\end{lstlisting}

\begin{lstlisting}[language=Python, caption={Construtor de \texttt{RobotSimulator}: compilação \texttt{lambdify} + CSE e inicialização do executor.}, label={lst:sim_init}]
from concurrent.futures import ProcessPoolExecutor
import time, numpy as np, sympy as sp

class RobotSimulator:
    def __init__(self, robot_math_instance, mode="Air"):
        self.bot       = robot_math_instance
        self.mode      = mode
        self.num_dof   = len(self.bot.q)
        self.q_home    = np.zeros(self.num_dof)
        self._executor = None      # executor persistente (inicializado apos compilacao)
        self._executor_workers = 0

        print(f"[{mode}] Compilando equacoes com lambdify+CSE...")
        _t0 = time.perf_counter()

        # Identificacao de juntas rotacionais (para wrap em [-pi, pi])
        self.is_rotational = np.array(
            [j.startswith('R') for j in self.bot.joint_config], dtype=bool)

        # Assinatura completa: [q1..qN, dq1..dqN, params...]
        self.sym_vars = self.bot.q + self.bot.dq + self.bot.params_list
        if hasattr(self.bot, 'rho'):
            self.sym_vars.append(self.bot.rho)

        # Compilacao separada por grandeza (M, C, G, J, FK)
        # cse=True: CSE antes da geracao do codigo -> reducao de 60-90% no bytecode
        self.expr_M = self.bot.M
        self.expr_C = self.bot.C_total
        self.expr_G = self.bot.G_vec

        self.func_M = sp.lambdify(self.sym_vars, self.expr_M,
                                  modules="numpy", cse=True)
        self.func_C = sp.lambdify(self.sym_vars, self.expr_C,
                                  modules="numpy", cse=True)
        self.func_G = sp.lambdify(self.sym_vars, self.expr_G,
                                  modules="numpy", cse=True)
        self.func_J = sp.lambdify(self.sym_vars, self.bot.Jacobian,
                                  modules="numpy", cse=True)

        # Funcoes FK por elo para visualizacao animada
        self.funcs_fk_all_links = [
            sp.lambdify(self.sym_vars,
                        frame[0:3, 3].subs(self.bot.step_4_prepare_export()["Subs"]),
                        modules="numpy", cse=True)
            for frame in self.bot.frames
        ]
        print(f"Compilacao concluida em {time.perf_counter()-_t0:.1f}s")

    def warm_up_workers(self, num_workers):
        """Inicializa o executor persistente e aquece os workers."""
        if (self._executor is not None
                and self._executor_workers == num_workers):
            return   # ja inicializado com mesmo numero de workers
        if self._executor is not None:
            self._executor.shutdown(wait=False)
        self._executor = ProcessPoolExecutor(
            max_workers=num_workers,
            initializer=_init_worker,
            initargs=(self.sym_vars, self.expr_M, self.expr_C, self.expr_G))
        self._executor_workers = num_workers
        # Warmup: submete tarefas dummy para pre-compilar os workers
        dummy_args = tuple(np.zeros(len(self.sym_vars)))
        futs = [self._executor.submit(_eval_worker, (k, dummy_args))
                for k in ("M", "C", "G")]
        for f in futs:
            try: f.result()
            except Exception: pass

    def close(self):
        """Encerra o executor persistente de forma limpa."""
        if self._executor is not None:
            self._executor.shutdown(wait=False)
            self._executor = None
\end{lstlisting}

% ==============================================================================
\chapter{Requisitos de Software, Instalação e Execução do \textit{Hephaestus}}
\label{ap:instalacao}

Este apêndice descreve os requisitos de software, o procedimento de instalação e as instruções de execução do ambiente \textit{Hephaestus} a partir do código-fonte.

% ------------------------------------------------------------------------------
\section*{C.1 Requisitos de Sistema}
% ------------------------------------------------------------------------------

O \textit{Hephaestus} foi desenvolvido e validado no seguinte ambiente:

\begin{itemize}
    \item \textbf{Sistema Operacional:} Windows 10/11 (64 bits). O código é multiplataforma e compatível com Linux e macOS com as mesmas dependências; no entanto, a interface gráfica com \textit{CustomTkinter} pode apresentar diferenças visuais menores entre plataformas.
    \item \textbf{Python:} versão 3.10 ou superior (recomendado: 3.11). O uso de \texttt{ProcessPoolExecutor} com método de \textit{spawn} requer Python $\geq$ 3.8; a tipagem e as dependências utilizadas requerem Python $\geq$ 3.10.
    \item \textbf{Hardware recomendado:} processador com pelo menos 4 núcleos lógicos (para aproveitar o paralelismo de derivação); 8 GB de RAM (para modelos com $N \geq 6$ GDL, cujas matrizes simbólicas podem ocupar centenas de MB antes da compilação).
\end{itemize}

% ------------------------------------------------------------------------------
\section*{C.2 Dependências Python}
% ------------------------------------------------------------------------------

A Tabela~\ref{tab:deps} lista as dependências diretas do projeto com as versões mínimas recomendadas.

\begin{table}[h!]
\centering
\caption{Dependências Python do \textit{Hephaestus}.}
\label{tab:deps}
\begin{tabular}{l l p{7cm}}
\hline
\textbf{Pacote} & \textbf{Versão mínima} & \textbf{Função no projeto} \\
\hline
\texttt{sympy}         & 1.12   & Álgebra simbólica, derivação de $M$, $C$, $G$ e \texttt{lambdify} \\
\texttt{numpy}         & 1.24   & Álgebra linear numérica, avaliação das funções compiladas \\
\texttt{matplotlib}    & 3.7    & Visualização dos gráficos de erro, torque e trajetória \\
\texttt{customtkinter} & 5.2    & Interface gráfica (\textit{GUI}) com tema moderno \\
\texttt{scipy}         & 1.11   & Funções auxiliares (interpolação, transformadas) \\
\hline
\end{tabular}
\end{table}

Todas as dependências listadas são pacotes de código aberto disponíveis no repositório PyPI.

% ------------------------------------------------------------------------------
\section*{C.3 Instalação}
% ------------------------------------------------------------------------------

Recomenda-se a criação de um ambiente virtual Python isolado para evitar conflitos entre versões de pacotes:

\begin{lstlisting}[style=bash, caption={Criação de ambiente virtual e instalação das dependências.}, label={lst:install}]
# Clonar o repositorio (ou extrair o arquivo comprimido)
git clone https://github.com/seu_usuario/hephaestus.git
cd hephaestus

# Criar e ativar ambiente virtual (Windows)
python -m venv .venv
.venv\Scripts\activate

# Instalar dependencias
pip install sympy numpy matplotlib customtkinter scipy

# Verificar instalacao
python -c "import sympy, numpy, matplotlib, customtkinter; print('OK')"
\end{lstlisting}

Para sistemas Linux/macOS, substituir \texttt{.venv\textbackslash Scripts\textbackslash activate} por \texttt{source .venv/bin/activate}.

% ------------------------------------------------------------------------------
\section*{C.4 Execução}
% ------------------------------------------------------------------------------

O ponto de entrada principal da aplicação é o arquivo \texttt{main.py}. Para iniciar a interface gráfica:

\begin{lstlisting}[style=bash, caption={Execução da interface gráfica do Hephaestus.}, label={lst:run}]
# Com o ambiente virtual ativado
python main.py
\end{lstlisting}

\textbf{Importante (Windows):} o \textit{Hephaestus} utiliza \texttt{ProcessPoolExecutor} com método de \textit{spawn} para o paralelismo de processos. Em sistemas Windows, é obrigatório que o ponto de entrada esteja protegido pelo bloco \texttt{if \_\_name\_\_ == '\_\_main\_\_':}, conforme implementado em \texttt{main.py}:

\begin{lstlisting}[language=Python, caption={Bloco de proteção obrigatório no \texttt{main.py} para Windows.}, label={lst:main_guard}]
import multiprocessing

if __name__ == '__main__':
    multiprocessing.freeze_support()   # necessario para executaveis PyInstaller
    app = App()
    app.mainloop()
\end{lstlisting}

A ausência deste bloco causaria a abertura recursiva de múltiplas janelas ao criar processos filhos no Windows.

% ------------------------------------------------------------------------------
\section*{C.5 Geração do Executável Autônomo (\textit{.exe})}
% ------------------------------------------------------------------------------

Para distribuir o \textit{Hephaestus} como um executável autônomo sem necessidade de instalação de Python, utiliza-se o \textit{PyInstaller}:

\begin{lstlisting}[style=bash, caption={Geração do executável PyInstaller.}, label={lst:pyinstaller}]
pip install pyinstaller

pyinstaller --onefile --windowed \
            --add-data "tooltip_utils.py;." \
            --name "Hephaestus" main.py
\end{lstlisting}

O arquivo resultante \texttt{dist/Hephaestus.exe} é executável sem instalação prévia de Python. Nesse modo, o suporte a \textit{spawn} de processos no Windows é detectado automaticamente pela condição \texttt{sys.frozen} do PyInstaller: quando o executável está \textit{frozen}, o número de \textit{workers} do \texttt{ProcessPoolExecutor} é automaticamente definido como 1 (modo serial), evitando o conflito de \textit{spawn} em módulos congelados. A derivação simbólica e a simulação operam normalmente, apenas sem paralelismo de processos.

% ------------------------------------------------------------------------------
\section*{C.6 Estrutura de Arquivos do Projeto}
% ------------------------------------------------------------------------------

\begin{table}[h!]
\centering
\caption{Estrutura principal de arquivos do repositório \textit{Hephaestus}.}
\label{tab:estrutura}
\begin{tabular}{l p{9.5cm}}
\hline
\textbf{Arquivo} & \textbf{Descrição} \\
\hline
\texttt{main.py}          & Ponto de entrada; interface gráfica (\textit{CustomTkinter}) com as abas Modelagem, Simulação e Análise \\
\texttt{engine.py}        & \textit{Engine} matemática: classes \texttt{RobotMathEngine} e \texttt{RobotMathHydro}; derivação simbólica de $M$, $C$, $G$, $J$ \\
\texttt{simulator.py}     & Simulador numérico: classe \texttt{RobotSimulator}; controladores \texttt{SlidingModeControl}, \texttt{SuperTwistingSMC}, \texttt{Decentralized\_LADRC}; integrador simplético \\
\texttt{tooltip\_utils.py} & Utilitários de interface: tooltips educacionais e blocos de ajuda contextual para os parâmetros físicos \\
\texttt{TCC/}             & Capítulos do documento TCC em \LaTeX{} \\
\hline
\end{tabular}
\end{table}

%
% Pacotes necessários no preâmbulo do documento principal:
%   \usepackage{listings}
%   \usepackage{xcolor}
% ==============================================================================

% Configuração global das listagens de código Python
\lstset{
  language=Python,
  basicstyle=\ttfamily\footnotesize,
  keywordstyle=\color{blue}\bfseries,
  commentstyle=\color{gray}\itshape,
  stringstyle=\color{teal},
  numberstyle=\tiny\color{gray},
  numbers=left,
  stepnumber=1,
  numbersep=8pt,
  backgroundcolor=\color{gray!6},
  frame=single,
  framerule=0.4pt,
  rulecolor=\color{gray!40},
  breaklines=true,
  breakatwhitespace=false,
  showstringspaces=false,
  tabsize=4,
  captionpos=b,
  morekeywords={self, True, False, None, yield, lambda, with, as, assert,
                del, elif, except, finally, from, global, import, in,
                is, not, or, pass, raise, try},
  literate={á}{{\'a}}1 {é}{{\'e}}1 {í}{{\'i}}1 {ó}{{\'o}}1 {ú}{{\'u}}1
           {â}{{\^a}}1 {ê}{{\^e}}1 {î}{{\^i}}1 {ô}{{\^o}}1 {û}{{\^u}}1
           {ã}{{\~a}}1 {õ}{{\~o}}1 {ç}{{\c{c}}}1
           {Á}{{\'A}}1 {É}{{\'E}}1 {Í}{{\'I}}1 {Ó}{{\'O}}1 {Ú}{{\'U}}1
           {Â}{{\^A}}1 {Ê}{{\^E}}1 {Ô}{{\^O}}1 {Ã}{{\~A}}1 {Õ}{{\~O}}1
           {Ç}{{\c{C}}}1,
}

\lstdefinestyle{bash}{
  language=bash,
  keywordstyle=\color{violet}\bfseries,
  commentstyle=\color{gray}\itshape,
  stringstyle=\color{teal},
  backgroundcolor=\color{gray!8},
}

\chapter{Código-fonte da \textit{Engine} Matemática (\texttt{engine.py})}
\label{ap:engine}

Este apêndice apresenta o código-fonte completo dos componentes fundamentais da \textit{engine} matemática do \textit{Hephaestus}, implementados no módulo \texttt{engine.py}. As listagens cobrem as funções de nível de módulo utilizadas para paralelismo, a classe \texttt{RobotMathEngine} (ambiente seco) e a subclasse \texttt{RobotMathHydro} (ambiente subaquático). O código é apresentado na íntegra para viabilizar reprodutibilidade e auditoria matemática das equações derivadas.

% ------------------------------------------------------------------------------
\section*{A.1 Funções de Nível de Módulo para Paralelismo}
% ------------------------------------------------------------------------------

As duas funções a seguir são definidas no escopo do módulo para permitir sua serialização via \texttt{pickle}, requisito do \texttt{ProcessPoolExecutor} no Python. A função \texttt{\_diff\_M\_col} diferencia a matriz de inércia em relação a uma variável de junta $q_k$; a função \texttt{\_calc\_coriolis\_row} monta a $i$-ésima linha do vetor de Coriolis pela fórmula dos símbolos de Christoffel.

\begin{lstlisting}[language=Python, caption={Funções paralelas para derivação de $\partial M/\partial q_k$ e cálculo de $C_i$.}, label={lst:parallel_funcs}]
from concurrent.futures import ProcessPoolExecutor
from itertools import repeat
import os, time
import sympy as sp
from sympy.physics.mechanics import dynamicsymbols

def _diff_M_col(args):
    """Diferencia a matriz M em relação a uma variável simbólica qk.
    Executado em processo separado para paralelizar o cálculo de dM/dq.
    """
    M_list, n, qk = args
    M = sp.Matrix(n, n, M_list)
    return list(M.diff(qk))


def _calc_coriolis_row(i, q, dq, dM_dq):
    n = len(q)
    termo_linha = sp.Integer(0)
    for j in range(n):
        for k in range(j, n):
            dM_ij_dk = dM_dq[k][i, j]
            dM_ik_dj = dM_dq[j][i, k]
            dM_jk_di = dM_dq[i][j, k]
            c_ijk = sp.Rational(1, 2) * (dM_ij_dk + dM_ik_dj - dM_jk_di)
            if c_ijk != 0:
                termo = c_ijk * dq[j] * dq[k]
                if k != j:
                    termo *= 2
                termo_linha += termo
    # sp.collect() omitido intencionalmente: custo superlinear sem benefício,
    # pois o cse=True no lambdify posterior realiza compactação sistematica.
    return termo_linha
\end{lstlisting}

% ------------------------------------------------------------------------------
\section*{A.2 Classe \texttt{RobotMathEngine} — Ambiente Seco}
% ------------------------------------------------------------------------------

\begin{lstlisting}[language=Python, caption={Construtor e configuração de símbolos de \texttt{RobotMathEngine}.}, label={lst:engine_init}]
class RobotMathEngine:
    def __init__(self, joint_config, link_vectors_mask,
                 logger_callback=None, num_workers=None):
        self.joint_config    = joint_config
        self.link_vectors_mask = [sp.Matrix(v) for v in link_vectors_mask]
        self.log             = logger_callback if logger_callback else print
        self.num_workers     = num_workers

        self.t = sp.symbols('t')
        self.g = sp.symbols('g')

        self.q, self.dq, self.params_list   = [], [], []
        self.frames, self.rotation_matrices  = [], []
        self.com_positions_global            = []
        self.angular_velocities              = []
        self.masses                          = []

        self.params_list.append(self.g)   # g sempre presente como parâmetro
        self.M, self.G_vec, self.C_total, self.Jacobian = None, None, None, None

    def _rot_matrix_local(self, axis, angle):
        c, s = sp.cos(angle), sp.sin(angle)
        if axis == 'x': return sp.Matrix([[1,0,0],[0,c,-s],[0,s,c]])
        if axis == 'y': return sp.Matrix([[c,0,s],[0,1,0],[-s,0,c]])
        if axis == 'z': return sp.Matrix([[c,-s,0],[s,c,0],[0,0,1]])
        return sp.eye(3)

    def _get_axis_vector(self, axis):
        if axis == 'x': return sp.Matrix([1, 0, 0])
        if axis == 'y': return sp.Matrix([0, 1, 0])
        if axis == 'z': return sp.Matrix([0, 0, 1])
        return sp.Matrix([0, 0, 0])
\end{lstlisting}

\begin{lstlisting}[language=Python, caption={Cinemática simbólica direta: transformações homogêneas e posições de CM.}, label={lst:step1}]
    def step_1_kinematics(self):
        self.log("1. (AR) Calculando Cinemarica (Com CM Arbitrario)...")
        T_acc   = sp.eye(4)
        R_acc   = sp.eye(3)
        omega_acc = sp.Matrix([0, 0, 0])

        for i, (j_type, link_vec) in enumerate(
                zip(self.joint_config, self.link_vectors_mask)):
            q  = dynamicsymbols(f'q{i+1}')
            dq = dynamicsymbols(f'q{i+1}', 1)
            self.q.append(q);  self.dq.append(dq)

            m  = sp.symbols(f'm{i+1}')
            L  = sp.symbols(f'L{i+1}')
            cx, cy, cz = sp.symbols(f'cx{i+1} cy{i+1} cz{i+1}')

            self.masses.append(m)
            self.params_list.extend([m, L, cx, cy, cz])

            type_char = j_type[0]
            axis_char = j_type[1].lower()

            if type_char == 'R':            # Junta de revolucao
                axis_vec_global = R_acc * self._get_axis_vector(axis_char)
                omega_new = omega_acc + axis_vec_global * dq
                R_j = self._rot_matrix_local(axis_char, q)
                P_j = sp.Matrix([0, 0, 0])
            elif type_char == 'D':          # Junta prismatica
                omega_new = omega_acc
                R_j = sp.eye(3)
                P_j = sp.Matrix([0, 0, 0])
                if axis_char == 'x': P_j[0] = q
                if axis_char == 'y': P_j[1] = q
                if axis_char == 'z': P_j[2] = q

            R_acc = R_acc * R_j
            self.rotation_matrices.append(R_acc)
            self.angular_velocities.append(omega_new)
            omega_acc = omega_new

            T_joint = sp.eye(4)
            T_joint[0:3, 0:3] = R_j
            T_joint[0:3, 3]   = P_j
            T_at_joint_start  = T_acc * T_joint

            P_link_vec        = link_vec * L
            T_link            = sp.eye(4)
            T_link[0:3, 3]    = P_link_vec
            T_acc             = T_at_joint_start * T_link
            self.frames.append(T_acc)

            # CM Global: transforma coordenadas locais pelo frame ate a junta
            v_cm_local   = sp.Matrix([cx, cy, cz, 1])
            p_cm_global  = T_at_joint_start * v_cm_local
            self.com_positions_global.append(p_cm_global[0:3, 0])
\end{lstlisting}

\begin{lstlisting}[language=Python, caption={Derivação simbólica de $M(q)$, $G(q)$ e $J(q)$ pelo método do Jacobiano da inércia.}, label={lst:step2}]
    def step_2_jacobian_M_G(self):
        self.log("2. (AR) Dinamica M e G (Steiner + Tensores)...")
        _t0 = time.perf_counter()
        n   = len(self.q)
        self.M = sp.zeros(n, n)
        V_tot  = 0

        for i in range(n):
            m       = self.masses[i]
            pos_cm  = self.com_positions_global[i]
            R_global = self.rotation_matrices[i]

            J_v = pos_cm.jacobian(self.q)         # Jacobiano linear do CM
            J_w = self.angular_velocities[i].jacobian(self.dq)  # Jacobiano angular

            Ixx, Iyy, Izz = sp.symbols(f'Ixx{i+1} Iyy{i+1} Izz{i+1}')
            self.params_list.extend([Ixx, Iyy, Izz])

            # Transformacao de Steiner: I_global = R * I_local * R^T
            I_local  = sp.Matrix([[Ixx, 0, 0], [0, Iyy, 0], [0, 0, Izz]])
            I_global = R_global * I_local * R_global.T

            # M += m*Jv^T*Jv + Jw^T * I_global * Jw  (metodo do Jacobiano)
            self.M  += m * J_v.T * J_v + J_w.T * I_global * J_w

            V_tot   += m * self.g * pos_cm[2]   # Energia potencial gravitacional
            self.log(f"  M/G: elo {i+1}/{n} concluido")

        self.G_vec   = sp.Matrix([V_tot]).jacobian(self.q).T
        J_lin        = self.frames[-1][0:3, 3].jacobian(self.q)
        J_ang        = self.angular_velocities[-1].jacobian(self.dq)
        self.Jacobian = J_lin.col_join(J_ang)    # Jacobiano 6xN completo
        self.log(f"  M/G/J: concluido em {time.perf_counter()-_t0:.1f}s")
\end{lstlisting}

\begin{lstlisting}[language=Python, caption={Cálculo do vetor de Coriolis pelos símbolos de Christoffel com paralelismo opcional.}, label={lst:step3}]
    def step_3_coriolis_combined(self):
        self.log("3. Calculando Coriolis...")
        n          = len(self.q)
        self.C_total = sp.zeros(n, 1)
        cpu_total  = os.cpu_count() or 1
        workers    = self.num_workers if self.num_workers is not None else cpu_total
        use_parallel = workers is not None and workers > 1

        # --- Passo 3a: dM/dq  (paralelizavel por coluna) ---
        self.log("  dM/dq: diferenciando M" +
                 (" em paralelo" if use_parallel else "") + "...")
        _t0 = time.perf_counter()
        M_list = list(self.M)
        tasks  = [(M_list, n, qk) for qk in self.q]

        if use_parallel:
            with ProcessPoolExecutor(max_workers=workers) as ex:
                dM_dq_flat = list(ex.map(_diff_M_col, tasks))
        else:
            dM_dq_flat = [_diff_M_col(t) for t in tasks]

        dM_dq = [sp.Matrix(n, n, flat) for flat in dM_dq_flat]
        self.log(f"  dM/dq: concluido em {time.perf_counter()-_t0:.1f}s")

        # --- Passo 3b: Simbolos de Christoffel  C_i(q,dq) ---
        if use_parallel:
            with ProcessPoolExecutor(max_workers=workers) as ex:
                for i, row in enumerate(
                    ex.map(_calc_coriolis_row,
                           range(n), repeat(self.q),
                           repeat(self.dq), repeat(dM_dq)),
                    start=1):
                    self.C_total[i - 1] = row
                    self.log(f"  Coriolis: linha {i}/{n} concluida")
        else:
            for i in range(n):
                self.C_total[i] = _calc_coriolis_row(i, self.q, self.dq, dM_dq)
                self.log(f"  Coriolis: linha {i+1}/{n} concluida")
\end{lstlisting}

\begin{lstlisting}[language=Python, caption={Substituição de símbolos dinâmicos e exportação do dicionário de modelo.}, label={lst:step4}]
    def step_4_prepare_export(self):
        self.log("4. Otimizando equacoes...")
        mapa_subs = {}
        for i in range(len(self.q)):
            mapa_subs[self.q[i]]  = sp.Symbol(f'q{i+1}')
            mapa_subs[self.dq[i]] = sp.Symbol(f'dq{i+1}')

        dq_vec      = sp.Matrix(self.dq)
        V_cartesian = self.Jacobian * dq_vec if self.Jacobian is not None else None

        return {
            "M":      self.M,
            "G":      self.G_vec,
            "C":      self.C_total,
            "J":      self.Jacobian,
            "FK_Pos": self.frames[-1],
            "FK_Vel": V_cartesian,
            "Subs":   mapa_subs,
            "Mode":   "Air"
        }
\end{lstlisting}

% ------------------------------------------------------------------------------
\section*{A.3 Subclasse \texttt{RobotMathHydro} — Extensão Hidrodinâmica}
% ------------------------------------------------------------------------------

\begin{lstlisting}[language=Python, caption={Extensão hidrodinâmica: massa adicionada, empuxo e potencial aparente.}, label={lst:hydro}]
class RobotMathHydro(RobotMathEngine):
    def __init__(self, joint_config, link_vectors_mask,
                 logger_callback=None, num_workers=None):
        super().__init__(joint_config, link_vectors_mask,
                         logger_callback, num_workers)
        self.rho     = sp.symbols('rho')   # densidade do fluido
        self.volumes = []

    def step_1_kinematics_hydro(self):
        super().step_1_kinematics()
        self.log("   -> Adicionando parametros hidrodinamicos...")
        for i in range(len(self.q)):
            vol = sp.symbols(f'vol{i+1}')
            self.volumes.append(vol)
            self.params_list.append(vol)   # vol_i adicionado a assinatura

    def step_2_jacobian_M_G(self):
        self.log("2. (AGUA) Dinamica: M (Inercia + Added Mass) e G (Peso - Empuxo)...")
        _t0  = time.perf_counter()
        n    = len(self.q)
        self.M = sp.zeros(n, n)
        V_tot  = 0

        for i in range(n):
            m       = self.masses[i]
            vol     = self.volumes[i]
            pos_cm  = self.com_positions_global[i]
            R_global = self.rotation_matrices[i]

            J_v = pos_cm.jacobian(self.q)
            J_w = self.angular_velocities[i].jacobian(self.dq)

            # Tensor de inercia rigido (referencial global via Steiner)
            Ixx, Iyy, Izz = sp.symbols(f'Ixx{i+1} Iyy{i+1} Izz{i+1}')
            self.params_list.extend([Ixx, Iyy, Izz])
            I_local_RB  = sp.Matrix([[Ixx,0,0],[0,Iyy,0],[0,0,Izz]])
            I_global_RB = R_global * I_local_RB * R_global.T

            # Tensor de massa adicionada (diagonal local, 6 componentes por elo)
            ma_u, ma_v, ma_w = sp.symbols(f'ma_u{i+1} ma_v{i+1} ma_w{i+1}')
            ma_p, ma_q, ma_r = sp.symbols(f'ma_p{i+1} ma_q{i+1} ma_r{i+1}')
            self.params_list.extend([ma_u, ma_v, ma_w, ma_p, ma_q, ma_r])

            MA_lin_local = sp.Matrix([[ma_u,0,0],[0,ma_v,0],[0,0,ma_w]])
            MA_rot_local = sp.Matrix([[ma_p,0,0],[0,ma_q,0],[0,0,ma_r]])

            # Rotacao para referencial global: M_A^(g) = R * M_A^(l) * R^T
            MA_lin_global = R_global * MA_lin_local * R_global.T
            MA_rot_global = R_global * MA_rot_local * R_global.T

            # M_hydro += Jv^T (m*I + M_A,lin)^(g) Jv + Jw^T (I_RB + M_A,rot)^(g) Jw
            self.M += J_v.T * (m * sp.eye(3) + MA_lin_global) * J_v
            self.M += J_w.T * (I_global_RB + MA_rot_global) * J_w

            # Potencial aparente: V_i = (m - rho*vol) * g * z_cm  (Peso - Empuxo)
            peso_aparente = (m - self.rho * vol) * self.g
            V_tot += peso_aparente * pos_cm[2]
            self.log(f"  M/G: elo {i+1}/{n} concluido")

        self.G_vec   = sp.Matrix([V_tot]).jacobian(self.q).T
        J_lin        = self.frames[-1][0:3, 3].jacobian(self.q)
        J_ang        = self.angular_velocities[-1].jacobian(self.dq)
        self.Jacobian = J_lin.col_join(J_ang)
        self.log(f"  M/G/J: concluido em {time.perf_counter()-_t0:.1f}s")

    def step_4_prepare_export(self):
        data = super().step_4_prepare_export()
        data["Mode"] = "Hydro"
        return data
\end{lstlisting}

% ==============================================================================
\chapter{Código-fonte dos Controladores e do Simulador (\texttt{simulator.py})}
\label{ap:simulador}

Este apêndice apresenta as classes de controle e a estrutura principal da classe \texttt{RobotSimulator} implementadas no módulo \texttt{simulator.py}. As listagens incluem os controladores CT-SMC, Super-Twisting (STA) e LADRC descentralizado, bem como o construtor do simulador, que realiza a compilação das expressões simbólicas via \texttt{lambdify} com CSE, e a infraestrutura de avaliação paralela numérica.

% ------------------------------------------------------------------------------
\section*{B.1 Funções Auxiliares de Orientação}
% ------------------------------------------------------------------------------

\begin{lstlisting}[language=Python, caption={Erro de orientação pelo operador de anel (skew-symmetric) e SLERP em $SO(3)$.}, label={lst:orient_helpers}]
import numpy as np

def _rotation_error_vec(R_ref, R_curr):
    """Erro de orientacao 3D pela formulacao de anel (skew-symmetric).

    Retorna e_R = vee(0.5*(R_ref @ R_curr^T - R_curr @ R_ref^T)),
    proporcional ao eixo-angulo para erros pequenos.
    """
    skew = 0.5 * (R_ref @ R_curr.T - R_curr @ R_ref.T)
    return np.array([skew[2, 1], skew[0, 2], skew[1, 0]])


def _slerp_rotation(R0, Rf, s, sd):
    """SLERP entre R0 e Rf na fracao s em [0, 1].

    Parâmetros
    ----------
    R0, Rf : matrizes de rotacao 3x3
    s      : fracao de interpolacao atual  (0 = R0, 1 = Rf)
    sd     : derivada temporal de s (ds/dt)

    Retorna
    -------
    R_ref     : rotacao interpolada 3x3
    omega_ref : velocidade angular no referencial global (rad/s)
    """
    dR        = R0.T @ Rf
    trace_val = np.clip((np.trace(dR) - 1.0) / 2.0, -1.0, 1.0)
    theta_total = np.arccos(trace_val)
    if abs(theta_total) < 1e-9:
        return R0.copy(), np.zeros(3)
    sin_t      = np.sin(theta_total)
    skew       = (dR - dR.T) / (2.0 * sin_t)
    axis_local = np.array([skew[2, 1], skew[0, 2], skew[1, 0]])
    theta_s    = s * theta_total
    K = np.array([[ 0.0,           -axis_local[2],  axis_local[1]],
                  [ axis_local[2],  0.0,            -axis_local[0]],
                  [-axis_local[1],  axis_local[0],   0.0          ]])
    dR_s  = np.eye(3) + np.sin(theta_s) * K + (1.0 - np.cos(theta_s)) * (K @ K)
    R_ref = R0 @ dR_s
    omega_ref = (sd * theta_total) * (R0 @ axis_local)
    return R_ref, omega_ref
\end{lstlisting}

% ------------------------------------------------------------------------------
\section*{B.2 Controlador CT-SMC com Camada Limite}
% ------------------------------------------------------------------------------

\begin{lstlisting}[language=Python, caption={Controle por Torque Computado com Modo Deslizante (CT-SMC).}, label={lst:smc}]
class SlidingModeControl:
    """Controlador CT-SMC: Computed Torque + Sliding Mode com camada limite.

    Lei de controle:
        q_ddot_f = alpha*q_ddot_d + (1-alpha)*q_ddot_f  (filtro passabaixa)
        v        = q_ddot_f - lambda*e_dot - K*tanh(S/phi)
        tau      = M(q)*v + C(q,dq) + G(q)
    onde S = e_dot + lambda*e  (superficie deslizante).

    Referencia: Slotine & Li, Applied Nonlinear Control, 1991.
    """

    def __init__(self, num_dof, lambda_gain, K, phi, ddq_filter_alpha=1.0):
        self.lambda_gain    = (np.full(num_dof, lambda_gain, dtype=float)
                               if np.isscalar(lambda_gain)
                               else np.array(lambda_gain, dtype=float))
        self.K              = (np.full(num_dof, K, dtype=float)
                               if np.isscalar(K) else np.array(K, dtype=float))
        self.phi            = (np.full(num_dof, phi, dtype=float)
                               if np.isscalar(phi) else np.array(phi, dtype=float))
        self.phi            = np.where(self.phi == 0, 1e-6, self.phi)
        self.alpha          = float(np.clip(ddq_filter_alpha, 1e-6, 1.0))
        self._ddq_d_filt    = None

    def compute_tau(self, q, dq, q_d, dq_d, ddq_d, M, C, G, dt=0.0):
        if self._ddq_d_filt is None:
            self._ddq_d_filt = ddq_d.copy()
        else:
            self._ddq_d_filt = (self.alpha * ddq_d
                                + (1.0 - self.alpha) * self._ddq_d_filt)
        e       = q - q_d
        e_dot   = dq - dq_d
        S       = e_dot + self.lambda_gain * e
        v       = (self._ddq_d_filt - self.lambda_gain * e_dot
                   - self.K * np.tanh(S / self.phi))
        return M @ v + C + G
\end{lstlisting}

% ------------------------------------------------------------------------------
\section*{B.3 Controlador Super-Twisting (STA)}
% ------------------------------------------------------------------------------

\begin{lstlisting}[language=Python, caption={Controlador Super-Twisting de 2ª ordem (STA) com sinal de controle contínuo.}, label={lst:sta}]
class SuperTwistingSMC:
    """Controlador Super-Twisting SMC (STA) baseado em modelo.

    Algoritmo de modo deslizante de 2a ordem:
        q_ddot_f = alpha*q_ddot_d + (1-alpha)*q_ddot_f
        v_sw     = -k1*|S|^(1/2)*sign(S) + z
        z_dot    = -k2*sign(S)
        tau      = M(q)*[q_ddot_f - lambda*e_dot + v_sw] + C(q,dq) + G(q)

    Propriedades: convergencia em tempo finito, sinal de controle continuo.
    Condicoes de ganho (Moreno & Osorio, 2012): k2 > delta, k1 > 2*sqrt(k2*delta).

    Referencias:
    - Levant, Int. J. Control, 58(6), 1993.
    - Moreno & Osorio, IEEE Trans. Autom. Control, 57(4), 2012.
    """

    def __init__(self, num_dof, lambda_gain, k1, k2, ddq_filter_alpha=1.0):
        self.lambda_gain = (np.full(num_dof, lambda_gain, dtype=float)
                            if np.isscalar(lambda_gain)
                            else np.array(lambda_gain, dtype=float))
        self.k1          = (np.full(num_dof, k1, dtype=float)
                            if np.isscalar(k1) else np.array(k1, dtype=float))
        self.k2          = (np.full(num_dof, k2, dtype=float)
                            if np.isscalar(k2) else np.array(k2, dtype=float))
        self.z           = np.zeros(num_dof)
        self.alpha       = float(np.clip(ddq_filter_alpha, 1e-6, 1.0))
        self._ddq_d_filt = None

    def reset(self):
        self.z[:]        = 0.0
        self._ddq_d_filt = None

    def compute_tau(self, q, dq, q_d, dq_d, ddq_d, M, C, G, dt):
        if self._ddq_d_filt is None:
            self._ddq_d_filt = ddq_d.copy()
        else:
            self._ddq_d_filt = (self.alpha * ddq_d
                                + (1.0 - self.alpha) * self._ddq_d_filt)
        e     = q - q_d
        e_dot = dq - dq_d
        S     = e_dot + self.lambda_gain * e

        # Parcela de chaveamento continua do STA
        v_sw = -self.k1 * np.abs(S) ** 0.5 * np.sign(S) + self.z

        v        = self._ddq_d_filt - self.lambda_gain * e_dot + v_sw
        self.z  -= self.k2 * np.sign(S) * dt    # integracao de z (Euler explicito)
        return M @ v + C + G
\end{lstlisting}

% ------------------------------------------------------------------------------
\section*{B.4 LADRC Descentralizado com ESO de Terceira Ordem}
% ------------------------------------------------------------------------------

\begin{lstlisting}[language=Python, caption={LADRC descentralizado: ESO de 3ª ordem e lei de controle por junta.}, label={lst:ladrc}]
class Decentralized_LADRC:
    """LADRC descentralizado para multiplas juntas (2a ordem por junta).

    Para a junta i: ddot_q_i = b0_i * tau_i + f_i  (f_i = disturbio total).
    O ESO de 3a ordem (linear) estima [q_i, dq_i, f_i] em z1, z2, z3.
    Parametrizacao dos ganhos (Gao, 2003): beta1=3*wo, beta2=3*wo^2, beta3=wo^3.
    """

    def __init__(self, num_dof, b0, wo, kp, kd, dt,
                 z_limit=None, tau_limit=None,
                 max_wo_dt=0.1, z3_filter_alpha=1.0):
        self.b0  = (np.full(num_dof, b0, dtype=float)
                    if np.isscalar(b0) else np.array(b0, dtype=float))
        self.wo  = (np.full(num_dof, wo, dtype=float)
                    if np.isscalar(wo) else np.array(wo, dtype=float))
        self.kp  = (np.full(num_dof, kp, dtype=float)
                    if np.isscalar(kp) else np.array(kp, dtype=float))
        self.kd  = (np.full(num_dof, kd, dtype=float)
                    if np.isscalar(kd) else np.array(kd, dtype=float))
        self.dt  = dt
        self.z_limit        = z_limit
        self.tau_limit      = tau_limit
        self.max_wo_dt      = max_wo_dt
        self.z3_filter_alpha = float(np.clip(z3_filter_alpha, 0.0, 1.0))
        self.z3_filtered    = None

        self._refresh_gains()
        self.z1 = np.zeros(num_dof, dtype=float)
        self.z2 = np.zeros(num_dof, dtype=float)
        self.z3 = np.zeros(num_dof, dtype=float)

    def _refresh_gains(self):
        # Limita wo para garantir estabilidade do ESO discreto: wo*dt <= max_wo_dt
        if self.max_wo_dt is not None:
            wo_dt = self.wo * self.dt
            self.wo = np.where(wo_dt > self.max_wo_dt,
                               self.max_wo_dt / self.dt, self.wo)
        self.beta1 = 3.0 *  self.wo
        self.beta2 = 3.0 * (self.wo ** 2)
        self.beta3 =         self.wo ** 3

    def reset_state(self, q, dq, z3=None):
        self.z1 = np.array(q,  dtype=float).copy()
        self.z2 = np.array(dq, dtype=float).copy()
        self.z3 = (np.zeros(self.num_dof, dtype=float) if z3 is None
                   else (np.full(self.num_dof, z3, dtype=float)
                         if np.isscalar(z3) else np.array(z3, dtype=float)))
        self.z3_filtered = self.z3.copy()

    def update_eso(self, q, tau_prev):
        """Avanca o ESO por um passo de integracao (Euler explicito)."""
        error  = self.z1 - q
        z1_dot = self.z2 - self.beta1 * error
        z2_dot = self.z3 - self.beta2 * error + self.b0 * tau_prev
        z3_dot =         - self.beta3 * error

        self.z1 += z1_dot * self.dt
        self.z2 += z2_dot * self.dt
        self.z3 += z3_dot * self.dt

        # Filtro IIR de 1a ordem em z3 (atenua alta frequencia antes da lei)
        a = self.z3_filter_alpha
        self.z3_filtered = a * self.z3 + (1.0 - a) * self.z3_filtered

    def compute_control(self, q_d, dq_d, ddq_d, q_meas=None, dq_meas=None):
        """Calcula o torque de controle LADRC (cancela disturbio + PD + ff)."""
        q_fb  = q_meas  if q_meas  is not None else self.z1
        dq_fb = dq_meas if dq_meas is not None else self.z2

        z3_use = self.z3_filtered if self.z3_filtered is not None else self.z3
        u0     = ddq_d + self.kp * (q_d - q_fb) + self.kd * (dq_d - dq_fb)
        b0_safe = np.where(self.b0 == 0.0, 1e-6, self.b0)
        u = (u0 - z3_use) / b0_safe
        if self.tau_limit is not None:
            u = np.clip(u, -self.tau_limit, self.tau_limit)
        return u
\end{lstlisting}

% ------------------------------------------------------------------------------
\section*{B.5 Compilação Simbólico-Numérica e Executor Persistente}
% ------------------------------------------------------------------------------

\begin{lstlisting}[language=Python, caption={Funções de worker para avaliação paralela de $M$, $C$, $G$ no executor persistente.}, label={lst:workers}]
import sympy as sp

_WORKER_FUNCS = {}   # dicionario de modulo: preenchido por _init_worker


def _init_worker(sym_vars, expr_M, expr_C, expr_G):
    """Inicializador do worker: compila lambdify localmente no processo."""
    global _WORKER_FUNCS
    _WORKER_FUNCS = {
        "M": sp.lambdify(sym_vars, expr_M, modules="numpy"),
        "C": sp.lambdify(sym_vars, expr_C, modules="numpy"),
        "G": sp.lambdify(sym_vars, expr_G, modules="numpy"),
    }


def _eval_worker(task):
    """Avalia uma das funcoes compiladas com os argumentos numericos fornecidos."""
    func_name, args = task
    return _WORKER_FUNCS[func_name](*args)
\end{lstlisting}

\begin{lstlisting}[language=Python, caption={Construtor de \texttt{RobotSimulator}: compilação \texttt{lambdify} + CSE e inicialização do executor.}, label={lst:sim_init}]
from concurrent.futures import ProcessPoolExecutor
import time, numpy as np, sympy as sp

class RobotSimulator:
    def __init__(self, robot_math_instance, mode="Air"):
        self.bot       = robot_math_instance
        self.mode      = mode
        self.num_dof   = len(self.bot.q)
        self.q_home    = np.zeros(self.num_dof)
        self._executor = None      # executor persistente (inicializado apos compilacao)
        self._executor_workers = 0

        print(f"[{mode}] Compilando equacoes com lambdify+CSE...")
        _t0 = time.perf_counter()

        # Identificacao de juntas rotacionais (para wrap em [-pi, pi])
        self.is_rotational = np.array(
            [j.startswith('R') for j in self.bot.joint_config], dtype=bool)

        # Assinatura completa: [q1..qN, dq1..dqN, params...]
        self.sym_vars = self.bot.q + self.bot.dq + self.bot.params_list
        if hasattr(self.bot, 'rho'):
            self.sym_vars.append(self.bot.rho)

        # Compilacao separada por grandeza (M, C, G, J, FK)
        # cse=True: CSE antes da geracao do codigo -> reducao de 60-90% no bytecode
        self.expr_M = self.bot.M
        self.expr_C = self.bot.C_total
        self.expr_G = self.bot.G_vec

        self.func_M = sp.lambdify(self.sym_vars, self.expr_M,
                                  modules="numpy", cse=True)
        self.func_C = sp.lambdify(self.sym_vars, self.expr_C,
                                  modules="numpy", cse=True)
        self.func_G = sp.lambdify(self.sym_vars, self.expr_G,
                                  modules="numpy", cse=True)
        self.func_J = sp.lambdify(self.sym_vars, self.bot.Jacobian,
                                  modules="numpy", cse=True)

        # Funcoes FK por elo para visualizacao animada
        self.funcs_fk_all_links = [
            sp.lambdify(self.sym_vars,
                        frame[0:3, 3].subs(self.bot.step_4_prepare_export()["Subs"]),
                        modules="numpy", cse=True)
            for frame in self.bot.frames
        ]
        print(f"Compilacao concluida em {time.perf_counter()-_t0:.1f}s")

    def warm_up_workers(self, num_workers):
        """Inicializa o executor persistente e aquece os workers."""
        if (self._executor is not None
                and self._executor_workers == num_workers):
            return   # ja inicializado com mesmo numero de workers
        if self._executor is not None:
            self._executor.shutdown(wait=False)
        self._executor = ProcessPoolExecutor(
            max_workers=num_workers,
            initializer=_init_worker,
            initargs=(self.sym_vars, self.expr_M, self.expr_C, self.expr_G))
        self._executor_workers = num_workers
        # Warmup: submete tarefas dummy para pre-compilar os workers
        dummy_args = tuple(np.zeros(len(self.sym_vars)))
        futs = [self._executor.submit(_eval_worker, (k, dummy_args))
                for k in ("M", "C", "G")]
        for f in futs:
            try: f.result()
            except Exception: pass

    def close(self):
        """Encerra o executor persistente de forma limpa."""
        if self._executor is not None:
            self._executor.shutdown(wait=False)
            self._executor = None
\end{lstlisting}

% ==============================================================================
\chapter{Requisitos de Software, Instalação e Execução do \textit{Hephaestus}}
\label{ap:instalacao}

Este apêndice descreve os requisitos de software, o procedimento de instalação e as instruções de execução do ambiente \textit{Hephaestus} a partir do código-fonte.

% ------------------------------------------------------------------------------
\section*{C.1 Requisitos de Sistema}
% ------------------------------------------------------------------------------

O \textit{Hephaestus} foi desenvolvido e validado no seguinte ambiente:

\begin{itemize}
    \item \textbf{Sistema Operacional:} Windows 10/11 (64 bits). O código é multiplataforma e compatível com Linux e macOS com as mesmas dependências; no entanto, a interface gráfica com \textit{CustomTkinter} pode apresentar diferenças visuais menores entre plataformas.
    \item \textbf{Python:} versão 3.10 ou superior (recomendado: 3.11). O uso de \texttt{ProcessPoolExecutor} com método de \textit{spawn} requer Python $\geq$ 3.8; a tipagem e as dependências utilizadas requerem Python $\geq$ 3.10.
    \item \textbf{Hardware recomendado:} processador com pelo menos 4 núcleos lógicos (para aproveitar o paralelismo de derivação); 8 GB de RAM (para modelos com $N \geq 6$ GDL, cujas matrizes simbólicas podem ocupar centenas de MB antes da compilação).
\end{itemize}

% ------------------------------------------------------------------------------
\section*{C.2 Dependências Python}
% ------------------------------------------------------------------------------

A Tabela~\ref{tab:deps} lista as dependências diretas do projeto com as versões mínimas recomendadas.

\begin{table}[h!]
\centering
\caption{Dependências Python do \textit{Hephaestus}.}
\label{tab:deps}
\begin{tabular}{l l p{7cm}}
\hline
\textbf{Pacote} & \textbf{Versão mínima} & \textbf{Função no projeto} \\
\hline
\texttt{sympy}         & 1.12   & Álgebra simbólica, derivação de $M$, $C$, $G$ e \texttt{lambdify} \\
\texttt{numpy}         & 1.24   & Álgebra linear numérica, avaliação das funções compiladas \\
\texttt{matplotlib}    & 3.7    & Visualização dos gráficos de erro, torque e trajetória \\
\texttt{customtkinter} & 5.2    & Interface gráfica (\textit{GUI}) com tema moderno \\
\texttt{scipy}         & 1.11   & Funções auxiliares (interpolação, transformadas) \\
\hline
\end{tabular}
\end{table}

Todas as dependências listadas são pacotes de código aberto disponíveis no repositório PyPI.

% ------------------------------------------------------------------------------
\section*{C.3 Instalação}
% ------------------------------------------------------------------------------

Recomenda-se a criação de um ambiente virtual Python isolado para evitar conflitos entre versões de pacotes:

\begin{lstlisting}[style=bash, caption={Criação de ambiente virtual e instalação das dependências.}, label={lst:install}]
# Clonar o repositorio (ou extrair o arquivo comprimido)
git clone https://github.com/seu_usuario/hephaestus.git
cd hephaestus

# Criar e ativar ambiente virtual (Windows)
python -m venv .venv
.venv\Scripts\activate

# Instalar dependencias
pip install sympy numpy matplotlib customtkinter scipy

# Verificar instalacao
python -c "import sympy, numpy, matplotlib, customtkinter; print('OK')"
\end{lstlisting}

Para sistemas Linux/macOS, substituir \texttt{.venv\textbackslash Scripts\textbackslash activate} por \texttt{source .venv/bin/activate}.

% ------------------------------------------------------------------------------
\section*{C.4 Execução}
% ------------------------------------------------------------------------------

O ponto de entrada principal da aplicação é o arquivo \texttt{main.py}. Para iniciar a interface gráfica:

\begin{lstlisting}[style=bash, caption={Execução da interface gráfica do Hephaestus.}, label={lst:run}]
# Com o ambiente virtual ativado
python main.py
\end{lstlisting}

\textbf{Importante (Windows):} o \textit{Hephaestus} utiliza \texttt{ProcessPoolExecutor} com método de \textit{spawn} para o paralelismo de processos. Em sistemas Windows, é obrigatório que o ponto de entrada esteja protegido pelo bloco \texttt{if \_\_name\_\_ == '\_\_main\_\_':}, conforme implementado em \texttt{main.py}:

\begin{lstlisting}[language=Python, caption={Bloco de proteção obrigatório no \texttt{main.py} para Windows.}, label={lst:main_guard}]
import multiprocessing

if __name__ == '__main__':
    multiprocessing.freeze_support()   # necessario para executaveis PyInstaller
    app = App()
    app.mainloop()
\end{lstlisting}

A ausência deste bloco causaria a abertura recursiva de múltiplas janelas ao criar processos filhos no Windows.

% ------------------------------------------------------------------------------
\section*{C.5 Geração do Executável Autônomo (\textit{.exe})}
% ------------------------------------------------------------------------------

Para distribuir o \textit{Hephaestus} como um executável autônomo sem necessidade de instalação de Python, utiliza-se o \textit{PyInstaller}:

\begin{lstlisting}[style=bash, caption={Geração do executável PyInstaller.}, label={lst:pyinstaller}]
pip install pyinstaller

pyinstaller --onefile --windowed \
            --add-data "tooltip_utils.py;." \
            --name "Hephaestus" main.py
\end{lstlisting}

O arquivo resultante \texttt{dist/Hephaestus.exe} é executável sem instalação prévia de Python. Nesse modo, o suporte a \textit{spawn} de processos no Windows é detectado automaticamente pela condição \texttt{sys.frozen} do PyInstaller: quando o executável está \textit{frozen}, o número de \textit{workers} do \texttt{ProcessPoolExecutor} é automaticamente definido como 1 (modo serial), evitando o conflito de \textit{spawn} em módulos congelados. A derivação simbólica e a simulação operam normalmente, apenas sem paralelismo de processos.

% ------------------------------------------------------------------------------
\section*{C.6 Estrutura de Arquivos do Projeto}
% ------------------------------------------------------------------------------

\begin{table}[h!]
\centering
\caption{Estrutura principal de arquivos do repositório \textit{Hephaestus}.}
\label{tab:estrutura}
\begin{tabular}{l p{9.5cm}}
\hline
\textbf{Arquivo} & \textbf{Descrição} \\
\hline
\texttt{main.py}          & Ponto de entrada; interface gráfica (\textit{CustomTkinter}) com as abas Modelagem, Simulação e Análise \\
\texttt{engine.py}        & \textit{Engine} matemática: classes \texttt{RobotMathEngine} e \texttt{RobotMathHydro}; derivação simbólica de $M$, $C$, $G$, $J$ \\
\texttt{simulator.py}     & Simulador numérico: classe \texttt{RobotSimulator}; controladores \texttt{SlidingModeControl}, \texttt{SuperTwistingSMC}, \texttt{Decentralized\_LADRC}; integrador simplético \\
\texttt{tooltip\_utils.py} & Utilitários de interface: tooltips educacionais e blocos de ajuda contextual para os parâmetros físicos \\
\texttt{TCC/}             & Capítulos do documento TCC em \LaTeX{} \\
\hline
\end{tabular}
\end{table}
